\BourbakiExerciseGroup

\begin{enumerate}
\def\labelenumi{\arabic{enumi})}
\tightlist
\item
  Let \(U = (\alpha_{ij})\) be a matrix with \(m\) rows and \(n\) columns over the principal ideal domain A. a) Suppose first of all that the \(\alpha_{ij}\) are setwise coprime. Show that there exist two invertible square matrices \(P\) and \(Q\) with entries in A, such that one of the entries of the matrix \(PUQ\) is equal to 1. (For each \(\alpha_{ij} \neq 0\) , let \(S(\alpha_{ij})\) be the sum of the exponents in a decomposition of \(\alpha_{ij}\) into irreducible elements, and let \(s(U)\) be the least of the numbers \(s(\alpha_{ij})\) ; if \(s(U) > 0\) , show that there exist two invertible square matrices \(R\) and \(S\) such that
\end{enumerate}

\[
s (R U S) <   s (U);
\]

use Bezout's identity, and restrict the matrices \(R\) and \(S\) considered, after permuting the rows and columns if necessary, to those of the form \(\begin{pmatrix} T & 0 \\ 0 & I \end{pmatrix}\) , where \(T\) is a \(2 \times 2\) matrix and \(I\) an identity matrix, or to products of such matrices and permutation matrices.)

\begin{enumerate}
\def\labelenumi{\alph{enumi})}
\setcounter{enumi}{1}
\tightlist
\item
  If \(\delta_1\) is a gcd of the entries of \(U\) , show that there exist two invertible matrices \(P_1, Q_1\) , such that
\end{enumerate}

\[
P _ {1} U Q _ {1} = \left( \begin{array}{c c} \delta_ {1} & 0 \\ 0 & U _ {1} \end{array} \right)
\]

where all the entries of \(U_{1}\) are divisible by \(\delta_{1}\) (use a)).

\begin{enumerate}
\def\labelenumi{\alph{enumi})}
\setcounter{enumi}{2}
\tightlist
\item
  For \(\mathbf{A} = \mathbf{Z}\) , use \(b\) to obtain a method of calculating the invariant factors of an explicit matrix with entries in \(\mathbf{Z}\) . Apply this method to the matrix
\end{enumerate}

\[
\left( \begin{array}{c c c c} 6 & 8 & 4 & 20 \\ 12 & 12 & 18 & 30 \\ 18 & 4 & 4 & 10 \end{array} \right).
\]

\begin{enumerate}
\def\labelenumi{\arabic{enumi})}
\setcounter{enumi}{1}
\item
  Let A be a principal ideal domain. Then two submodules M and N of \(A^{q}\) are transforms of one another under automorphisms of \(A^{q}\) if and only if they have the same invariant factors with respect to \(A^{q}\) .
\item
  Let A be a principal ideal domain with field of fractions K, let E be a vector space over K, let M be a finitely generated A-module of rank n contained in E, and let N be a finitely generated A-module of rank \(p \leqslant n\) contained in KM. Show that there exists a basis \((e_{i})_{1 \leqslant i \leqslant n}\) of M and a basis of N consisting of vectors \(\alpha_{i}e_{i}\) ( \(1 \leqslant i \leqslant p\) ), where \(\alpha_{i} \in K\) and \(\alpha_{i}\) divides \(\alpha_{i+1}\) (with respect to the ring A) for \(1 \leqslant i \leqslant p - 1\) ; show also that the fractional ideals \(A\alpha_{i}\) are uniquely determined (invariant factors of N with respect to M).
\item
  Let \(G\) be a finite abelian group.
\end{enumerate}

\begin{enumerate}
\def\labelenumi{\alph{enumi})}
\item
  For every prime number \(p\) , show that the number of elements of order \(p\) in \(G\) is \(p^N - 1\) , where \(N = \sum_{n=1}^{\infty} m(p^n)\) (in the notation of Prop. 9 of VII, p.~24).
\item
  Deduce that if the equation \(x^p = 1\) has at most \(p\) solutions in \(G\) , for each prime number \(p\) , then \(G\) is cyclic.
\end{enumerate}

\begin{enumerate}
\def\labelenumi{\arabic{enumi})}
\setcounter{enumi}{4}
\item
  Let \(G\) be a finite abelian group of order \(n\) . Show that for every integer \(q\) dividing \(n\) there exists a subgroup of order \(q\) in \(G\) (decompose \(G\) as a direct sum of indecomposable cyclic subgroups).
\item
  Let A be a commutative ring and let E be an A-module. Let \(\mathcal{E}\) denote the ring \(\operatorname{End}_{\mathbf{z}}(\mathbf{E})\) of endomorphisms of the abelian group E (without operators); for every subring \(\mathbf{B} \subset \mathcal{E}\) , let \(\mathbf{B}'\) denote the centraliser of B in \(\mathcal{E}\) (the subring of \(\mathcal{E}\) consisting of elements which commute with every element of B).
\end{enumerate}

\begin{enumerate}
\def\labelenumi{\alph{enumi})}
\item
  Let \(A_{0}\) be the subring of E consisting of the homotheties \(x \mapsto \lambda x (\lambda \in A)\) ; if a is the annihilator of E, then \(A_{0}\) is known to be isomorphic to A/a, and its centraliser \(A_{0}^{\prime}\) in E is the ring \(\operatorname{End}_{\mathbf{A}}(\mathbf{E})\) of endomorphisms of E as an A-module. Show that the centraliser \(A_{0}^{\prime\prime}\) of \(A_{0}^{\prime}\) is commutative, and that its centraliser \(A_{0}^{\prime\prime\prime}\) is equal to \(A_{0}^{\prime}\) (notice that \(A_{0} \subset A_{0}^{\prime}\) ).
\item
  If the submodule F is a direct factor of E, show that \(u(F) \subset F\) for all \(u \in A_{0}''\) (consider the projection of E onto F obtained from the decomposition of E as a direct sum of F and another submodule).
\item
  Suppose that E is the direct sum of a family \((F_{i})\) of cyclic submodules. If \(u \in A_{0}''\) show that, for each index i, there exists an element \(\alpha_{i} \in A\) such that \(u(x) = \alpha_{i}x\) for all \(x \in F_{i}\) (use b)).
\item
  Suppose E is the direct sum of a sequence \((\mathrm{E}_{n})\) (finite or infinite) of cyclic submodules such that, if \(b_{n}\) denotes the annihilator of \(E_{n}\) , then the sequence \((\mathfrak{b}_{n})\) is increasing. Show that \(A_{0}^{\prime\prime}=A_{0}\) (use c), and the fact that for each index i\textgreater1 there exists an endomorphism \(v_{i}\in A_{0}^{\prime}\) which maps \(E_{i-1}\) into \(E_{i}\) .
\end{enumerate}

Apply this to the case where A is a principal ideal domain and E is a finitely generated A-module.

\begin{enumerate}
\def\labelenumi{\arabic{enumi})}
\setcounter{enumi}{6}
\tightlist
\item
  Let A be a principal ideal domain with field of fractions K, and let E be a finitely generated, torsion-free A-module of rank n.~Let \(E^{*}\) be the dual of E; then E and \(E^{*}\) are each isomorphic to \(A^{n}\) .
\end{enumerate}

\begin{enumerate}
\def\labelenumi{\alph{enumi})}
\item
  Let M be a submodule of E; show that, if \(M^{0}\) is the submodule of \(E^{*}\) orthogonal to M, then the module \(E^{*}/M^{0}\) is torsion-free, and consequently \(M^{0}\) is a direct factor in \(E^{*}\) ; the submodule \(M^{00}\) of E orthogonal to \(M^{0}\) is equal to \(E \cap KM\) (considering E as being canonically embedded in an n-dimensional vector space over K). We can consider \(E^{*}/M^{0}\) as being canonically embedded in the dual \(M^{*}\) of M; show that the invariant factors of \(E^{*}/M^{0}\) with respect to \(M^{*}\) are equal to the invariant factors of M with respect to E (use Th. 1 of VII, p.~18).
\item
  Let F be a second finitely generated, torsion-free A-module, and let u be a linear transformation from E into F. Show that the invariant factors of \(u(F^{*})\) with respect to \(E^{*}\) are equal to those of \(u(E)\) with respect to F (proceed as in Prop. 5 of VII, p.~21).
\end{enumerate}

\begin{enumerate}
\def\labelenumi{\arabic{enumi})}
\setcounter{enumi}{7}
\tightlist
\item
  Let G be a module of finite length over a principal ideal domain A. For any pair of free modules M and N such that \(N \subset M\) and G is isomorphic to M/N, the invariant factors of N with respect to M which are distinct from A are identical to the invariant factors of G (VII, p.~20, Def. 2); their number is called the rank of G.
\end{enumerate}

\begin{enumerate}
\def\labelenumi{\alph{enumi})}
\item
  Show that the rank r of G is the least number of cyclic submodules of G of which G is the direct sum, and is equal to the greatest of the ranks of the \(\pi\) -primary components of G. b) Show that the rank of a submodule or a quotient module of G is at most equal to the rank of G (consider G as a quotient of two free modules of rank r).
\item
  For each \(\lambda \in \mathbf{A}\) , show that the rank of the submodule \(\lambda G\) is equal to the number of invariant factors of \(G\) which do not divide \(\lambda\) (notice that, if \(E = A / (A\alpha)\) , then the module \(\lambda E\) is isomorphic to \((A\lambda) / ((A\alpha) \cap (A\lambda))\) . Deduce that the \(k\) -th invariant factor of \(G\) (in decreasing order) is a gcd of those \(\lambda \in A\) such that \(\lambda G\) has rank \(\leqslant r - k\) .
\item
  Let \(A\alpha_{k}\) ( \(1 \leqslant k \leqslant r\) ) be the invariant factors of G arranged in decreasing order. Let H be a submodule of G of rank r - q, and let \(A\beta_{k}\) ( \(1 \leqslant k \leqslant r - q\) ) be its invariant factors, arranged in decreasing order. Deduce from b) and c) that \(\beta_{k}\) divides \(\alpha_{k+q}\) for \(1 \leqslant k \leqslant r - q\) . Similarly show that, if G/H has rank r - p and invariant factors \(A\gamma_{k}\) ( \(1 \leqslant k \leqslant r - p\) ) in decreasing order, then \(\gamma_{k}\) divides \(\alpha_{k+p}\) for \(1 \leqslant k \leqslant r - p\) .
\item
  Conversely, let L be a torsion module of rank r - q, whose invariant factors \(A\lambda_{k}\) ( \(1 \leqslant k \leqslant r - q\) ) are such that \(\lambda_{k}\) divides \(\alpha_{k+q}\) for \(1 \leqslant k \leqslant r - q\) . Show that there exist submodules M and N of G such that L is isomorphic to M and to G/N (decompose G into a direct sum of submodules isomorphic to A/A \(\alpha_{k}\) ).
\end{enumerate}

\begin{enumerate}
\def\labelenumi{\arabic{enumi})}
\setcounter{enumi}{8}
\tightlist
\item
  Let A be a principal ideal domain and K its field of fractions; let E be a vector space over K, let M be a finitely generated A-module of rank n contained in E, let N be a finitely generated A-module of rank p contained in KM, and let P be a finitely generated A-module of rank q contained in KN.
\end{enumerate}

\begin{enumerate}
\def\labelenumi{\alph{enumi})}
\item
  Let \(\mathbf{A}\alpha_{i}\) ( \(1 \leqslant i \leqslant p\) ) be the invariant factors of \(\mathbf{N}\) with respect to \(\mathbf{M}\) , arranged in decreasing order (Ex. 3). Show that, for \(1 \leqslant k \leqslant p\) , the element \(\alpha_{k}\) is a gcd of those elements \(\lambda \in \mathbf{K}\) such that the rank of the quotient module \((\mathbf{N} + \lambda (\mathbf{KN} \cap \mathbf{M})) / \mathbf{N}\) is \(\leqslant p - k\) (cf.~Ex. 8, d)).
\item
  Let \(\mathbf{A}\beta_{j}\) ( \(1 \leqslant j \leqslant q\) ) be the invariant factors of \(\mathbf{P}\) with respect to \(\mathbf{M}\) and \(\mathbf{A}\gamma_{j}\) ( \(1 \leqslant j \leqslant q\) ) the invariant factors of \(\mathbf{P}\) with respect to \(\mathbf{N}\) , arranged in decreasing order. Show that \(\alpha_{1}\gamma_{j}\) divides \(\beta_{j}\) for \(1 \leqslant j \leqslant q\) (cf.~Ex. 8).
\item
  Show that \(\gamma_{1}\alpha_{j}\) divides \(\beta_{j}\) for \(1 \leqslant j \leqslant q\) . (First consider the case \(n = p = q\) , and apply Ex. 8. In general, show that we may always assume \(p = n\) , and consider a submodule Q of M such that \(\mathbf{P} + \rho \mathbf{Q}\) has rank \(p\) for all \(\rho \in \mathbf{A}\) ; then choose the ideal \(\mathbf{A}\rho\) sufficiently small, and apply Prop. 4 of VII, p.~20.)
\end{enumerate}

\(d\) ) For any submodule H of N of rank \(k \leqslant p\) , let \((\mu_{\mathrm{H}})\) be the fractional ideal consisting of \(\mu \in \mathbf{K}\) such that \(\mu (\mathbf{M} \cap \mathbf{KH}) \subset \mathbf{H}\) . Deduce from \(c\) ) that for each \(1 \leqslant k \leqslant p\) the element \(\alpha_{k}\) is a gcd of the \(\mu_{\mathrm{H}}\) as H runs through all submodules of N of rank \(k\) .

\begin{enumerate}
\def\labelenumi{\arabic{enumi})}
\setcounter{enumi}{9}
\tightlist
\item
  Let A be a principal ideal domain and let M be a submodule of rank n in \(E = A^{n}\) ; let \(A\alpha_{i} (1 \leqslant i \leqslant n)\) be the invariant factors of M with respect to E, arranged in decreasing order. Let N be a submodule of M admitting a complement in M (that is, such that \(N = M \cap KN\) , where K is the field of fractions of A).
\end{enumerate}

\begin{enumerate}
\def\labelenumi{\alph{enumi})}
\item
  Let P be a complement of N in M. Show that the submodule \((P + (E \cap KN))/M\) of E/M is isomorphic to \((E \cap KN)/N\) . Deduce that, if N has rank p and \(A\beta_{k}\) ( \(1 \leqslant k \leqslant p\) ) are the invariant factors of N with respect to E, then \(\beta_{k}\) divides \(\alpha_{k+n-p}\) for \(1 \leqslant k \leqslant p\) (use Ex. 8, d)).
\item
  Give an example to show that the module E/M is not necessarily isomorphic to the product of the modules \((E \cap KN)/N\) and \(E/(P + (E \cap KN))\) (take E/M to be indecomposable).
\end{enumerate}

\begin{enumerate}
\def\labelenumi{\arabic{enumi})}
\setcounter{enumi}{10}
\item
  Let A be a principal ideal domain, let H be a submodule of \(E = A^{n}\) , and let \(H^{0}\) be the submodule of the dual \(E^{*}\) of E orthogonal to H (Ex. 7). Let M be a submodule of rank p of E and let \(A\alpha_{i}\) ( \(1 \leqslant i \leqslant p\) ) be the invariant factors of M with respect to E, arranged in decreasing order. Let \((\nu_{\mathrm{H}})\) be the gcd of the ideals \(f(\mathbf{M})\) of A as f runs through \(H^{0}\) ; show that, for \(1 \leqslant k \leqslant p\) , the element \(\alpha_{k}\) is a gcd of the \(\nu_{H}\) for all submodules H of E of rank k - 1 which admit a complement in E (first consider the case k = 1; apply the result to the submodule \((\mathbf{M} + \mathbf{H})/\mathbf{H}\) of E/H, and use Ex. 8, d), applied to the quotient module \(E/(M + H)\) of E/M).
\item
  Let G be a module of finite length over a ring A. Show that a submodule H of G can only be isomorphic to G if it is equal to G (cf.~II, p.~212, Prop. 16).
\item
  Let A be a principal ideal domain with field of fractions K.
\end{enumerate}

\begin{enumerate}
\def\labelenumi{\alph{enumi})}
\item
  Show that, for every A-module M, the canonical homomorphism \(c_{\mathbf{M}}: \mathbf{M} \to \mathbf{D}(\mathbf{D}(\mathbf{M}))\) is injective (cf.~II, p.~389, Ex. 13).
\item
  If \(u: \mathbf{M} \to \mathbf{N}\) is a homomorphism of A-modules, define natural isomorphisms from \(D(\operatorname{Im}(u))\) onto \(\operatorname{Im}(D(u))\) , from \(D(\operatorname{Ker}(u))\) onto \(\operatorname{Coker}(D(u))\) , and from \(D(\operatorname{Coker}(u))\) onto \(\operatorname{Ker}(D(u))\) .
\item
  Let M be an A-module and N a submodule of M. In the notation of Sect. 9, show that \(N^{00} = N\) (use a) and b)). The module \(N^{0}\) has finite length if and only if M/N has finite length. Conversely, for any submodule \(N'\) of D(M) of finite length, the module \(N'^{0}\) is a submodule of M such that \(M/N^{0}\) has finite length, and \(N' = N'^{00}\) .
\end{enumerate}

\(d\) ) If \(\mathbf{M}\) is a torsion A-module, and \(\mathbf{M}_{\pi}\) denotes the \(\pi\) -primary component of \(\mathbf{M}\) , show that \(\mathrm{D}(\mathbf{M})\) is isomorphic to \(\prod_{\pi \in \mathbb{P}}\mathrm{Hom}(\mathbf{M}_{\pi},\mathbf{U}_{\pi})\) , where \(\mathbf{P}\) is a system of representatives of irreducible elements of A and \(U_{\pi}\) is the \(\pi\) -primary component of K/A for each \(\pi \in P\) (VII, p.~54 Ex. 3).

\begin{enumerate}
\def\labelenumi{\alph{enumi})}
\setcounter{enumi}{4}
\tightlist
\item
  Deduce from d) that \(\mathrm{D}(\mathrm{K}/\mathrm{A})\) is isomorphic to the product \(\prod_{\pi\in P}\hat{\mathbf{A}}_{(\pi)}\) , where
\end{enumerate}

\(\hat{\mathbf{A}}_{(\pi)}\) denotes the completion of the ring \(\mathbf{A}_{(\pi)}\) defined in VII, p.~54, Ex. 2, with respect to the topology defined by taking the ideals \(\pi^n\mathbf{A}_{(\pi)}\) to be a basis of open neighbourhoods of 0 (this completion can also be identified with the inverse limit \(\lim (\mathbf{A} / \mathbf{A}\pi^n)\) ).

\begin{enumerate}
\def\labelenumi{\alph{enumi})}
\setcounter{enumi}{5}
\item
  If M is an A-module of finite length, and a is its annihilator, show that D(M) is isomorphic to the dual of the faithful (A/a)-module associated to M.
\item
  Let M be an A-module of finite length. Show that if N is a submodule of M, then there exists a submodule P of M such that M/N is isomorphic to P and M/P is isomorphic to N (use VII, p.~26, Prop. 10 and VII, p.~27, Prop. 12); two submodules of M with these properties are called reciprocal. Show that, for all \(\alpha \in A\) , the submodule \(M(\alpha)\) of M consisting of all \(x \in M\) such that \(\alpha x = 0\) , and the submodule \(\alpha M\) , are reciprocal.
\end{enumerate}

¶ 14) Let M be a module of finite length over a principal ideal domain A with field of fractions K.

\begin{enumerate}
\def\labelenumi{\alph{enumi})}
\tightlist
\item
  Show that if N is a submodule of M then the rank of M is at most equal to the sum of the ranks of N and M/N (let M = E/H, where E is a free module of rank n and H is a submodule of rank n of E; if N = L/H, where H ⊆ L ⊆ E, notice that, if the rank of N is p, then there exists a submodule R of H, of rank n - p, admitting a complement S in L such that S ∩ H is a complement of R in H; moreover, (E ∩ KR)/R has rank ≥ n - p and is isomorphic to a quotient module of E/L (Ex. 10, a)); complete the proof using Ex. 8, b)). b) Let N be a submodule of M of rank p, and let q be the rank of M/N. Let Aα \(_{i}\) (1 ≤ i ≤ n) be the invariant factors of M and Aβ \(_{k}\) (1 ≤ k ≤ p) those of N, arranged in decreasing order. Show that β \(_{k}\) is a multiple of α \(_{k- (p + q - n)}\) (use a) and Ex. 8, c), noticing that (λM)/(λN) is isomorphic to a quotient module of M/N).
\end{enumerate}

\begin{enumerate}
\def\labelenumi{\arabic{enumi})}
\setcounter{enumi}{14}
\tightlist
\item
  \begin{enumerate}
  \def\labelenumii{\alph{enumii})}
  \tightlist
  \item
    Let M be a module of finite length and rank n over a principal ideal domain A, and let \(\mathbf{A}\alpha_{i}\) ( \(1 \leqslant i \leqslant n\) ) be its invariant factors in decreasing order. Let \((\beta_{i})\) ( \(1 \leqslant i \leqslant n\) ) be a sequence of elements of A such that \(\beta_{i}\) divides \(\alpha_{i}\) for \(1 \leqslant i \leqslant n\) and \(\alpha_{i}\) divides \(\beta_{i+1}\) for \(1 \leqslant i \leqslant n-1\) . Let \((\mathbf{M}_{i})_{1 \leqslant i \leqslant n}\) be a sequence of submodules of M such that M is the direct sum of the \(M_{i}\) and each \(M_{i}\) is isomorphic to \(A/A\alpha_{i}\) ; let \(a_{i}\) be a generator of \(M_{i}\) for \(1 \leqslant i \leqslant n\) . Put
  \end{enumerate}
\end{enumerate}

\[
b = \beta_ {1} a _ {1} + \frac {\beta_ {1} \beta_ {2}}{\alpha_ {1}} a _ {2} + \frac {\beta_ {1} \beta_ {2} \beta_ {3}}{\alpha_ {1} \alpha_ {2}} a _ {3} + \dots + \frac {\beta_ {1} \beta_ {2} \dots \beta_ {n}}{\alpha_ {1} \alpha_ {2} \dots \alpha_ {n - 1}} a _ {n}.
\]

Show that the quotient of M by the cyclic submodule Ab is isomorphic to the direct sum of the n modules \(A/A\beta_{i}\) ( \(1 \leq i \leq n\) ).

\begin{enumerate}
\def\labelenumi{\alph{enumi})}
\setcounter{enumi}{1}
\tightlist
\item
  Let M and N be two modules of finite length over A; let \(A\alpha_{i}\) ( \(1 \leqslant i \leqslant n\) ) and \(A\beta_{k}\) ( \(1 \leqslant k \leqslant p\) ) be the invariant factors of M and N respectively, arranged in decreasing order; suppose that \(p \leqslant n\) , that \(\beta_{k}\) divides \(\alpha_{k+n-p}\) for \(1 \leqslant k \leqslant p\) , and that \(\alpha_{k}\) divides \(\beta_{k+(p+q-n)}\) for \(1 \leqslant k \leqslant n-q\) , where q is an integer such that
\end{enumerate}

\[
q \leqslant n \leqslant p + q.
\]

Show that there exists a submodule P of M, isomorphic to N, such that M/P has rank \(\leq q\) (decompose each of the two modules M and N as a direct sum of q modules, in such a way as to reduce the problem to the case q = 1; then use a) and Ex. 13, g)).

\begin{enumerate}
\def\labelenumi{\alph{enumi})}
\setcounter{enumi}{2}
\tightlist
\item
  For a module N over A to be isomorphic to a direct factor of a module M of finite length, it is necessary and sufficient that every elementary divisor of N be an elementary divisor of M, and that its multiplicity in N be at most equal to its multiplicity in M. Hence obtain an example of a module M of finite length and a submodule N of M such that the rank of M is equal to the sum of the ranks of N and M/N, and such that N does not admit a complement in M (use a)).
\end{enumerate}

\begin{enumerate}
\def\labelenumi{\arabic{enumi})}
\setcounter{enumi}{15}
\tightlist
\item
  Let M be a module over a commutative ring A. A submodule N of M is called characteristic if \(u(N) \subset N\) for every endomorphism u of M. For all \(\alpha \in A\) , the submodule \(M(\alpha)\) of elements annihilated by \(\alpha\) , and the submodule \(\alpha M\) , are characteristic.
\end{enumerate}

\begin{enumerate}
\def\labelenumi{\alph{enumi})}
\item
  If N is a characteristic submodule of M, and P and Q are two complementary submodules of M, show that N is the direct sum of \(N \cap P\) and \(N \cap Q\) (consider the projections from M onto P and Q).
\item
  Let M be a module of finite length over a principal ideal domain A; let \(A\alpha_{i}\) be the invariant factors of M, arranged in decreasing order \((1 \leqslant i \leqslant n)\) ; let \((\mathbf{M}_{i})_{1 \leqslant i \leqslant n}\) be a sequence of submodules of M such that M is the direct sum of the \(M_{i}\) and each \(M_{i}\) is isomorphic to \(A/A\alpha_{i}\) . Let N be a characteristic submodule of M and let \(A\beta_{i}\) be the annihilator of \(N \cap M_{i}\) \((1 \leqslant i \leqslant n)\) ; if \(\alpha_{i} = \beta_{i}\gamma_{i}\) , show that \(\beta_{i}\) divides \(\beta_{i+1}\) and \(\gamma_{i}\) divides \(\gamma_{i+1}\) for all \(1 \leqslant i \leqslant n-1\) (notice that, for \(1 \leqslant i \leqslant n-1\) , there exists an endomorphism of M which maps \(M_{i+1}\) onto \(M_{i}\) , and an endomorphism of M which maps \(M_{i}\) onto a submodule of \(M_{i+1}\) ).
\item
  Conversely, let N be a submodule of M which is a direct sum of submodules \(N_{i} \subset M_{i}\) ; suppose the annihilators \(A\beta_{i}\) of the \(N_{i}\) ( \(1 \leqslant i \leqslant n\) ) satisfy the above conditions. Then show that N is a characteristic submodule of M. Deduce that if two characteristic submodules of M have the same invariant factors then they are equal.
\end{enumerate}

\begin{enumerate}
\def\labelenumi{\arabic{enumi})}
\setcounter{enumi}{16}
\tightlist
\item
  Let A be a principal ideal domain, let \(\alpha\) be a nonzero element of A, and let E be the module A/A \(\alpha\) . Then the product \(E^{n}\) is isomorphic to \(A^{n}/(\alpha A^{n})\) .
\end{enumerate}

\begin{enumerate}
\def\labelenumi{\alph{enumi})}
\item
  Let \(u\) be a linear map from \(\mathbf{E}^n\) into \(\mathbf{E}^m\) ; show that \(u\) can be obtained from a linear map \(v\) from \(\mathbf{A}^n\) into \(\mathbf{A}^m\) on passing to quotients.
\item
  Let \(A\beta_{i}\) ( \(1 \leqslant i \leqslant p \leqslant \inf(m, n)\) ) be the invariant factors of \(v(A^{n})\) with respect to \(A^{m}\) , arranged in decreasing order; let \(q \leqslant p\) be the greatest of the indices i such that \(\alpha\) does not divide \(\beta_{i}\) , and let \(\delta_{i}\) be a gcd of \(\alpha\) and \(\beta_{i}\) for \(1 \leqslant i \leqslant q\) . Show that the kernel \(u^{-1}(0)\) of u is the direct sum of n - q modules isomorphic to E and q cyclic modules isomorphic to \(A/A\delta_{i}\) respectively ( \(1 \leqslant i \leqslant q\) ). If \(\alpha = \gamma_{i}\delta_{i}\) ( \(1 \leqslant i \leqslant q\) ), show that the module \(u(E^{n})\) has invariant factors \(A\gamma_{i}\) .
\end{enumerate}

\begin{enumerate}
\def\labelenumi{\arabic{enumi})}
\setcounter{enumi}{17}
\tightlist
\item
  Let \(\mathbf{C}\) be a principal ideal domain. Consider a system of linear equations
\end{enumerate}

\[
\sum_ {j = 1} ^ {n} \alpha_ {i j} \xi_ {j} = \beta_ {i} (1 \leqslant i \leqslant m)
\]

where the coefficients and the right hand sides belong to C, and let A denote the matrix \((\alpha_{ij})\) with m rows and n columns, and B the matrix obtained by extending A by the \((n+1)\) -st column \((\beta_{i})\) . Show that the system admits at least one solution in \(C^{n}\) if and only if the following conditions are satisfied: 1) A and B have the same rank p; 2) a gcd of the minors of order p of A is equal to a gcd of the minors of order p of B (use Prop. 4 of VII, p.~20 and Ex. 9, c)).

\begin{enumerate}
\def\labelenumi{\arabic{enumi})}
\setcounter{enumi}{18}
\tightlist
\item
  Let \(\mathbf{C}\) be a principal ideal domain. Consider a system of \(m\) « linear congruences » \(\sum_{j=1}^{n} \alpha_{ij} \xi_j \equiv \beta_i (\bmod \alpha)\) , where the coefficients, the right hand sides, and the element \(\alpha\) all belong to C, and where \(\alpha\) is nonzero; let \(A\) denote the matrix \((\alpha_{ij})\) , let \(B\) denote the matrix obtained by extending A by the column \((\beta_i)\) , and let \(\delta_k\) (resp. \(\delta_k'\) ) be a gcd of the \(k\) -th order minors of \(A\) (resp. \(B\) ). Then the system admits at least one solution in \(\mathbf{C}^n\) if and only if: \(a)\) If \(m \leqslant n\) , a gcd of \(\alpha^m\) , \(\alpha^{m-1}\delta_1\) , \ldots, \(\alpha\delta_{m-1}\) , \(\delta_m\) is equal to a gcd of \(\alpha^m\) , \(\alpha^{m-1}\delta_1'\) , \ldots, \(\alpha\delta_{m-1}'\) , \(\delta_m'\) ;
\end{enumerate}

\begin{enumerate}
\def\labelenumi{\alph{enumi})}
\setcounter{enumi}{1}
\tightlist
\item
  if \(m > n\) , a gcd of \(\alpha^{n+1}\) , \(\alpha^n \delta_1, \ldots, \alpha \delta_n\) is equal to a gcd of \(\alpha^{n+1}\) , \(\alpha^n \delta_1', \ldots, \delta_{n+1}'\) .
\end{enumerate}

(Reduce the system of linear congruences to a system of linear equations, and use Ex. 18.)

\begin{enumerate}
\def\labelenumi{\arabic{enumi})}
\setcounter{enumi}{19}
\tightlist
\item
  \begin{enumerate}
  \def\labelenumii{\alph{enumii})}
  \tightlist
  \item
    Let \(\mathbf{C}\) be a principal ideal domain, and let
  \end{enumerate}
\end{enumerate}

\[
f (x _ {1},.., x _ {p}) = \sum_ {(i _ {k})} \alpha_ {i _ {1} \dots i _ {p}} \xi_ {1, i _ {1}} \dots \xi_ {p, i _ {p}}
\]

be a \(p\) -linear form on \(\mathbf{E}^p\) , where \(\mathbf{E}\) is the C-module \(\mathbf{C}^n\) (and where \(\xi_{ij}\) denotes the \(j\) -th coordinate of \(x_i\) ). Show that if \(\delta\) is a gcd of the coefficients \(\alpha_{i_1\dots i_p}\) then there exists an element \((a_k)\in \mathbf{E}^p\) such that \(f(a_1,\dots,a_p) = \delta\) (reduce to the case \(\delta = 1\) , and argue by induction on \(p\) , using Th. 1 of VII, p.~18).

\begin{enumerate}
\def\labelenumi{\alph{enumi})}
\setcounter{enumi}{1}
\tightlist
\item
  Let n \textgreater{} 1 be an integer, and \(\alpha_{i}\) ( \(1 \leqslant i \leqslant n\) ) be n elements of C with gcd \(\delta\) . Show that there exists a square matrix A of order n over C whose first column is \((\alpha_{i})\) and whose determinant is \(\delta\) (use a)).
\end{enumerate}

\begin{enumerate}
\def\labelenumi{\arabic{enumi})}
\setcounter{enumi}{20}
\tightlist
\item
  Let \(C\) be a principal ideal domain, and let \(A = (\alpha_{ij})\) be a square matrix of order \(n\) over \(C\) .
\end{enumerate}

\begin{enumerate}
\def\labelenumi{\alph{enumi})}
\item
  Show that there exists a square matrix U of order n over C, with determinant 1, such that the last column of the matrix UA has its first \((n-1)\) entries zero, and its last entry equal to the gcd of the entries in the n-th column of A (use Ex. 20, b)).
\item
  Deduce from a) that there exists a square matrix V of order n over C, with determinant 1, such that the matrix \(VA = (\beta_{ij})\) is lower triangular (« Hermite reduced form of A »).
\item
  If W is a square matrix of order n over C, with determinant 1, such that \(WA = (\gamma_{ij})\) is lower triangular, show that \(\beta_{ii}\) and \(\gamma_{ii}\) are associate for \(1 \leqslant i \leqslant n\) .
\end{enumerate}

\begin{enumerate}
\def\labelenumi{\arabic{enumi})}
\setcounter{enumi}{21}
\tightlist
\item
  A module M over a principal ideal domain is said to have finite rank if there exists an integer \(m \geqslant 0\) with the following property: every finitely generated submodule of M is contained in a submodule generated by m elements or fewer. The smallest integer m with this property is called the rank of M.
\end{enumerate}

\begin{enumerate}
\def\labelenumi{\alph{enumi})}
\item
  Show that if \(\mathbf{M}\) has finite rank, then every quotient module has finite rank, at most equal to that of \(\mathbf{M}\) .
\item
  For a torsion-free module M, the notion of rank defined above coincides with that defined in II, p.~315.
\item
  Show that if \(\mathbf{M}\) has finite rank then every submodule of \(\mathbf{M}\) has finite rank, at most equal to that of \(\mathbf{M}\) .
\item
  Let \(\mathbf{M}\) be a \(\pi\) -primary module such that \(\mathbf{M}(\pi) = \mathbf{M}\) (in the notation of VII, p.~7). Show that if \(\mathbf{M}\) has finite rank then it is finitely generated.
\item
  Let M be a \(\pi\) -primary module with no element of infinite height. Show that if M has finite rank then it is finitely generated, and the notion of rank coincides with that defined in
\end{enumerate}

Ex. 8. (Notice that if \(\mathbf{M}(\pi)\) is finitely generated then so is \(\mathbf{M}\) , by proving that the heights of elements of \(\mathbf{M}(\pi)\) in \(\mathbf{M}\) are bounded; use Ex. 5 of VII, p.~55 to show this.) \(f)\) Let \(\mathbf{M}\) be a \(\pi\) -primary module of finite rank \(r\) , and let \(\mathbf{M}_0\) denote the submodule consisting of all the elements of infinite height in \(\mathbf{M}\) . Show that \(\mathbf{M}_0\) is a divisible module, and the direct sum of a finite number \(p \leqslant r\) of indecomposable submodules \(\mathrm{U}_{\pi}\) (VII, p.~54, Ex. 3), and that \(\mathbf{M}\) is the direct sum of \(\mathbf{M}_0\) and a finitely generated submodule \(\mathbf{N}\) of rank \(r - p\) . (Notice that \(\mathbf{M} / \mathbf{M}_0\) is finitely generated, using \(e\) ); then apply Ex. 3 of VII, p.~54.) \(g)\) Deduce from the above that an A-module \(\mathbf{M}\) has finite rank \(r\) if and only if the quotient module \(\mathbf{M} / \mathbf{T}\) of \(\mathbf{M}\) by its torsion submodule \(\mathbf{T}\) has finite rank \(p \leqslant r\) , and that each \(\pi\) -primary component of \(\mathbf{T}\) has rank \(\leqslant r - p\) , with at least one of them having rank \(r - p\) .

\begin{enumerate}
\def\labelenumi{\alph{enumi})}
\setcounter{enumi}{7}
\item
  Let \((p_{n})\) be the sequence of distinct prime numbers; let E be the direct sum of the cyclic Z-modules \(Z/p_{n}^{2}Z\) , and let \(e_{n}\) denote the element whose coordinates all vanish except the nth, which is equal to the residue class of \(1 \mod p_{n}^{2}\) . Let M be the submodule of the Z-module \(E \times Q\) generated by the elements \((e_{n}, 1/p_{n})\) . Show that M is a Z-module of rank 2, that the torsion submodule T of M is generated by the elements \((p_{n}e_{n}, 0)\) , and that T is not a direct factor of M. (If \(\bar{x}_{n}\) is the class of \((e_{n}, 1/p_{n}) \mod T\) , show that for all \(m \neq n\) there exists a class \(\bar{y}_{nm} \mod T\) such that \(\bar{x}_{n} = p_{m}\bar{y}_{nm}\) , but that there is no element \(x_{n}\) in the class \(\bar{x}_{n}\) such that for all \(m \neq n\) there exists \(y_{nm} \in M\) with \(x_{n} = p_{m}y_{nm}\) .)
\item
  If \(\mathbf{M}\) is a finitely generated A-module, show that the rank of \(\mathbf{M}\) is equal to the smallest cardinal \(\gamma (\mathbf{M})\) of any generating set for \(\mathbf{M}\) .
\item
  Let M be a torsion A-module of finite rank. Show that if N is a proper submodule of M then N cannot be isomorphic to M (use f).
\end{enumerate}

\begin{enumerate}
\def\labelenumi{\arabic{enumi})}
\setcounter{enumi}{22}
\item
  Let A be a principal ideal domain and let M be a finitely generated A-module. Show that if N and P are two A-modules such that \(M \oplus N\) and \(M \oplus P\) are isomorphic, then N and P are isomorphic. (Reduce to the case where M is cyclic; by identifying \(M \oplus N\) with \(M \oplus P\) , we may consider an A-module E decomposed in two ways as a direct sum of submodules \(F \oplus N\) and \(G \oplus P\) , where F and G are cyclic and isomorphic. In the case where F and G are isomorphic to A, show that if \(N \neq P\) then N is the direct sum of \(N \cap P\) and a submodule isomorphic to A. If F and G are isomorphic to \(A/A\pi^{n}\) , where \(\pi\) is an irreducible element of A, consider two generators a, b of F, G respectively, and examine the cases \(\pi^{n-1}a \notin G\) , \(\pi^{n-1}b \notin F\) , and finally the case where both \(\pi^{n-1}a \in G\) and \(\pi^{n-1}b \in F\) ; in this last case, consider the element \(a + b\) of E and the submodule it generates.)
\item
  An abelian group G is isomorphic to a subgroup of the multiplicative group \(K^{*}\) of nonzero elements of some commutative field K if and only if the torsion subgroup T of G has rank \(\leqslant 1\) (Ex. 22). (For the necessity of this condition, cf.~V, p.~79, Prop. 2; use the same reference to show that the condition is sufficient for T to be isomorphic to a subgroup of \(E^{*}\) , for some field E. For all \(\alpha \in G/T\) let \(g_{\alpha}\) be a coset representative of \(\alpha\) in G, then, identifying T with a subgroup of \(E^{*}\) , we may write \(g_{\alpha}g_{\beta} = e_{\alpha\beta}g_{(\alpha\beta)}\) for all \(\alpha\) , \(\beta \in G/T\) , where \(e_{\alpha\beta} \in E\) . Show that there exists an E-algebra B with basis \((b_{\alpha})_{\alpha \in G/T}\) such that \(b_{\alpha}b_{\beta} = e_{\alpha\beta}b_{(\alpha\beta)}\) for all \(\alpha\) , \(\beta\) , and prove that B is an integral domain using VI, p.~33, Ex. 20, c). Then the field of fractions K of B satisfies the desired conditions.)
\item
  Generalise Th. 1 of VII, p.~18 to left modules over a ring A which is not necessarily commutative, in which every left ideal and every right ideal is principal (an example of such a ring is the tensor product \(\mathbf{K}[\mathbf{X}] \otimes_{\mathbf{K}} \mathbf{D}\) , where \(\mathbf{K}\) is a commutative field and \(\mathbf{D}\) is a K-algebra which is a field but not necessarily commutative).
\end{enumerate}

\BourbakiExerciseGroup

All fields in this section are commutative unless explicitly stated otherwise.

\begin{enumerate}
\def\labelenumi{\arabic{enumi})}
\item
  Let u be an endomorphism of a finite dimensional vector space E, and let V be a subspace of E invariant under u. Show that the characteristic polynomial of the restriction of u to V divides that of u. Hence obtain an elementary proof of the Cayley-Hamilton theorem (reduce to the case where \(E_{u}\) is cyclic, and use formula (2) of VII, p.~30).
\item
  \begin{enumerate}
  \def\labelenumii{\alph{enumii})}
  \tightlist
  \item
    Show that every square matrix over a field K is similar to its transpose (reduce to the case of a Jordan matrix).
  \end{enumerate}
\end{enumerate}

\begin{enumerate}
\def\labelenumi{\alph{enumi})}
\setcounter{enumi}{1}
\item
  Show that over the ring \(\mathbf{Z}\) of rational integers, the matrix \(\begin{pmatrix} 8 & 2 \\ 0 & 1 \end{pmatrix}\) is not similar to its transpose.
\item
  Let \(u\) be an endomorphism of a finite dimensional vector space over \(\mathbf{K}\) . Then every eigenvalue \(\lambda\) of \(u\) is also an eigenvalue of the transpose endomorphism \(^t u\) of \(E^*\) . Moreover, if \(\lambda\) and \(\mu\) are two distinct eigenvalues of \(u\) , if \(x\) is an eigenvector of \(u\) corresponding to \(\lambda\) , and if \(y^*\) is an eigenvector of \(^t u\) corresponding to \(\mu\) , then \(\langle x, y^* \rangle = 0\) .
\end{enumerate}

\begin{enumerate}
\def\labelenumi{\arabic{enumi})}
\setcounter{enumi}{2}
\item
  Let u be an endomorphism of a finite dimensional vector space E, let \(\lambda\) be an eigenvalue of u, and let v be the endomorphism \(u - \lambda\) . Show that E is the direct sum of \(v^{-1}(0)\) and \(v(E)\) if and only if \(\lambda\) is a simple root of the minimal polynomial of u.
\item
  Show, by adjoining a transcendental element to the base field, that formula (6) in Prop. 10 of VII, p.~36 can be deduced from formula (8). Prove the latter directly by decomposing the polynomial q into linear factors.
\item
  \begin{enumerate}
  \def\labelenumii{\alph{enumii})}
  \tightlist
  \item
    Let K be an algebraically closed field, and let A be a square matrix of order n over K, in the form of a diagonal block \(\text{diag}(A_i)\) of lower triangular matrices \(A_i\) , such that the diagonal entries of \(A_i\) are all equal to the same element \(\alpha_i\) of K ( \(1 \leqslant i \leqslant r\) ), and \(\alpha_i \neq \alpha_j\) for \(i \neq j\) . Show that every square matrix B which commutes with A is a diagonal block \(\text{diag}(B_i)\) , where \(B_i\) is the same size as \(A_i\) for \(1 \leqslant i \leqslant r\) .
  \end{enumerate}
\end{enumerate}

\begin{enumerate}
\def\labelenumi{\alph{enumi})}
\setcounter{enumi}{1}
\tightlist
\item
  Let M be a set of commuting square matrices of order n over K. Show that by applying the same similarity to all the matrices in M, we can make them all into diagonal blocks \(\text{diag}(X_i)\) of lower triangular matrices, where the number r of \(X_i\) 's and the order \(m_i\) of \(X_i\) for each i, are the same for each matrix in M, and where all the eigenvalues of \(X_i\) are equal. (Use a) to reduce to the case where r = 1, then Lemma 3 of VII, p.~41 to show that in this case all the matrices in M have nonzero common eigenvector, and complete the proof by induction on n.)
\end{enumerate}

\begin{enumerate}
\def\labelenumi{\arabic{enumi})}
\setcounter{enumi}{5}
\item
  Two pairs \((A_{1}, A_{2})\) and \((B_{1}, B_{2})\) of square matrices of order n over a commutative field K are said to be equivalent if there exist two invertible square matrices P and Q over K such that \(A_{1} = PB_{1}Q\) and \(A_{2} = PB_{2}Q\) . If \(A_{1}\) and \(B_{1}\) are invertible, show that the pairs are equivalent if and only if the matrices \(XA_{1} + A_{2}\) and \(XB_{1} + B_{2}\) are equivalent over the ring K{[}X{]}.
\item
  Let \(A\) be a matrix with \(m\) rows and \(n\) columns, of rank \(r\) , over a field \(K\) . If \(P\) is a square matrix of order \(m\) and \(Q\) is a square matrix of order \(n\) over \(K\) , such that \(PAQ = A\) , show that there exists a matrix \(P'\) similar to \(P\) and a matrix \(Q'\) similar to \(Q\) such that
\end{enumerate}

\[
P ^ {\prime} = \left( \begin{array}{c c} R & S \\ 0 & U \end{array} \right), \quad Q ^ {\prime} = \left( \begin{array}{c c} R ^ {- 1} & 0 \\ S ^ {\prime} & U ^ {\prime} \end{array} \right)
\]

where R is an invertible square matrix of order r, U (resp. \(U'\) ) is a square matrix of order m - r (resp. n - r), and S (resp. \(S'\) ) is a matrix with r rows and m - r columns (resp. n - r rows and m columns).

\begin{enumerate}
\def\labelenumi{\arabic{enumi})}
\setcounter{enumi}{7}
\item
  Let E be a vector space of finite dimension n over a field K, let u be an endomorphism of E, let V be a subspace of E invariant under u, and let \(u|V\) be the restriction of u to V. If \(f_{i}\) ( \(1 \leqslant i \leqslant n\) ) (resp. \(g_{j}\) ( \(1 \leqslant j \leqslant m\) )) are the similarity invariants of u (resp. \(u|V\) ), show that \(g_{j}\) divides \(f_{j+n-m}\) for \(1 \leqslant j \leqslant m\) (cf.~VII, p.~64, Ex. 8, d)). State and prove the analogous property of the endomorphism of E/V induced by u.
\item
  Let \(A\) be a square matrix of order \(n\) , let \(B\) be a square matrix or order \(m\) , and let \(C\) be a matrix with \(m\) rows and \(n\) columns over a field \(K\) . Reduce the equation \(X A = B X + C\) , where \(X\) is an unknown matrix with \(m\) rows and \(n\) columns over \(K\) , to a system of linear congruences in the polynomial ring \(K[Y]\) (for all \(f \in K[Y]\) , notice that \(X f(A) = f(B)X + H(f)\) for some well defined matrix \(H(f)\) with \(m\) rows and \(n\) columns; let \(\oplus_{i} E_i\) be the decomposition of the module \(E\) associated to \(A\) (Sect. 1) as a direct sum of
\end{enumerate}

cyclic modules, whose annihilators are the similarity invariants \(f_{i}\) of A which are distinct from 1 (Prop. 2), and let \(a_{i}\) be generators of the \(E_{i}\) ; in the same way define the submodules \(F_{j}\) , generated by \(b_{j}\) , and the polynomials \(g_{j}\) corresponding to B; write \(Xf_{i}(A)a_{i}=0\) ; then decompose \(Xa_{i}\) and \(H(f_{i})a_{i}\) in terms of the \(F_{j}\) , and show that the conditions thus obtained from the equations \(Xf_{i}(A)a_{i}=0\) are necessary and sufficient for the existence of a solution).

\begin{enumerate}
\def\labelenumi{\arabic{enumi})}
\setcounter{enumi}{9}
\tightlist
\item
  \begin{enumerate}
  \def\labelenumii{\alph{enumii})}
  \tightlist
  \item
    Show that, over an arbitrary field K, the matrix
  \end{enumerate}
\end{enumerate}

\[
\left( \begin{array}{c c c c c c} \lambda & \alpha & \beta_ {13} & \beta_ {14} & ... & \beta_ {1n} \\ 0 & \lambda & \alpha & \beta_ {24} & ... & \beta_ {2n} \\ 0 & 0 & \lambda & \alpha & ... & \beta_ {3n} \\ \cdots & \cdots & \cdots & \cdots & \cdots & \cdots \\ 0 & 0 & 0 & 0 & ... & \alpha \\ 0 & 0 & 0 & 0 & ... & \lambda \end{array} \right)
\]

is similar to the Jordan matrix \(U_{n,\lambda}\) if \(\alpha \neq 0\) .

\begin{enumerate}
\def\labelenumi{\alph{enumi})}
\setcounter{enumi}{1}
\item
  If \(\lambda \neq 0\) and \(\mathbf{K}\) is of characteristic 0, show that \((U_{n,\lambda})^m\) is similar to the Jordan matrix \(U_{n,\lambda^m}\) for every integer \(m\) (positive or negative).
\item
  Show that \((U_{n,0})^m = 0\) for \(m \geqslant n\) . For \(0 < m < n\) , let \(n - 1 = km + q\) with \(0 \leqslant q < m\) ; show that \((U_{n,0})^m\) has \(q + 1\) similarity invariants equal to \(\mathbf{X}^{k+1}\) and \(m - q - 1\) similarity invariants equal to \(\mathbf{X}^k\) (arrange the canonical basis of \(\mathbf{K}^n\) in a suitable order).
\item
  Similarly, determine the similarity invariants of \(f(U_{n,\lambda})\) for \(f \in \mathbf{K}[\mathbf{X}]\) .
\item
  Discuss the same problems for K of characteristic \(p \neq 0\) .
\end{enumerate}

\begin{enumerate}
\def\labelenumi{\arabic{enumi})}
\setcounter{enumi}{10}
\item
  Let K be an algebraically closed field, let A be an invertible matrix over K, and m \textgreater{} 0 an integer. Show that, if m is not a multiple of the characteristic of K, then there exists a polynomial \(g(\mathbf{X}) \in \mathbf{K}[\mathbf{X}]\) such that \((g(A))^{m} = A\) (use the Cayley-Hamilton theorem and Ex. 28 of VII, p.~54). Generalise to the equation \(f(U) = A\) ( \(f \in K[X]\) , and U an unknown \(n \times n\) matrix over K).
\item
  \begin{enumerate}
  \def\labelenumii{\alph{enumii})}
  \tightlist
  \item
    Let \(A\) and \(B\) be two square matrices of order \(n\) over a commutative field \(\mathbf{K}\) , and let \(f_i\) and \(g_i\) ( \(1 \leqslant i \leqslant n\) ) be their respective similarity invariants. Show that the vector space (over \(\mathbf{K}\) ) of square matrices \(U\) , of order \(n\) , such that \(UA = BU\) , is isomorphic to the space \(\mathbf{M}_0 = \mathbf{M} / \mathbf{N}\) defined as follows: \(\mathbf{M}\) is the space of square matrices \((u_{ij}(\mathbf{X}))\) of order \(n\) over \(\mathbf{K}[\mathbf{X}]\) such that \(u_{ij}(\mathbf{X}) f_j(\mathbf{X})\) is a multiple of \(g_i(\mathbf{X})\) for each pair \(i, j\) of indices; and \(\mathbf{N}\) is the subspace of \(\mathbf{M}\) consisting of those matrices for which \(u_{ij}(\mathbf{X})\) is a multiple of \(g_i(\mathbf{X})\) for each pair \(i, j\) of indices. (Consider \(U, A\) and \(B\) as the matrices of endomorphisms \(u, v\) and \(w\) respectively of \(E = K^n\) , and notice that \(u\) is a \(K[X]\) -linear map from \(E_v\) into \(E_w\) ; now express \(E_v\) and \(E_w\) as quotient modules.)
  \end{enumerate}
\end{enumerate}

\begin{enumerate}
\def\labelenumi{\alph{enumi})}
\setcounter{enumi}{1}
\item
  When B = A, the matrices U such that UA = AU form a ring P; then M is a ring of square matrices, and N a 2-sided ideal of M, such that P is isomorphic to the quotient ring M/N. Show that P is identical to the ring K{[}A{]} generated by \(I_{n}\) and A if and only if the minimal polynomial of A is equal to its characteristic polynomial.
\item
  For arbitrary \(A\) and \(B\) , show that the dimension of \(\mathbf{M}_0\) is equal to \(\sum_{i,j} m_{ij}\) , where
\end{enumerate}

\(m_{ij}\) is the degree of the gcd of \(f_i\) and \(g_j\) (same method as \(a\) ).

\begin{enumerate}
\def\labelenumi{\alph{enumi})}
\setcounter{enumi}{3}
\tightlist
\item
  Show that the greatest of the ranks of the matrices \(U \in \mathbf{M}_0\) is equal to \(\sum_{k} d_k\) , where
\end{enumerate}

\(d_{k}\) is the degree of the gcd of \(f_{k}\) and \(\pmb{g}_{k}\) (use Ex. 8).

\begin{enumerate}
\def\labelenumi{\alph{enumi})}
\setcounter{enumi}{4}
\tightlist
\item
  Generalise the results of a), c) and d) to the case where A is a square matrix of order n, B is a square matrix of order m, and U is a matrix with m rows and n columns.
\end{enumerate}

\begin{enumerate}
\def\labelenumi{\arabic{enumi})}
\setcounter{enumi}{12}
\tightlist
\item
  Let \((\mathbf{Z}_{ij})\) ( \(1 \leqslant i, j \leqslant n\) ) be \(n^2\) indeterminates over a field \(\mathbf{K}_0\) . Let \(\mathbf{A}\) be the polynomial ring \(\mathbf{K}_0[\mathbf{Z}_{11}, \ldots, \mathbf{Z}_{nn}]\) , and let \(\mathbf{K} = \mathbf{K}_0(\mathbf{Z}_{ij})\) be its field of fractions; let \(\chi_U\) be the characteristic polynomial of the matrix \(U = (\mathbf{Z}_{ij})\) over \(\mathbf{K}\) .
\end{enumerate}

\begin{enumerate}
\def\labelenumi{\alph{enumi})}
\item
  Show that \(\chi_{U}\) is a separable irreducible polynomial (show that, if \(\chi_{U}\) were reducible over K, it would be the product of two polynomials in \(K_{0}[X, Z_{ij}]\) of degree \textgreater0 in X; hence deduce that the characteristic polynomial of every square matrix of order n over K would be reducible; then derive a contradiction from the example of a polynomial with algebraically independent coefficients over \(K_{0}\) ; use the analogous method to prove separability).
\item
  Show that U is similar to a matrix in which every entry is equal to 0, 1 or a coefficient of the polynomial \(\chi_{U}\) (cf.~VII, p.~30, formula (2)). Deduce that, if \(K_{0}\) is an infinite field, and p is a polynomial in A such that \(p(u_{ij}^{\prime}) = p(u_{ij}^{\prime\prime})\) whenever the matrices \((u_{ij}^{\prime})\) and \((u_{ij}^{\prime\prime})\) over \(K_{0}\) are similar, then p belongs to the subring of A generated by the coefficients of \(\chi_{U}\) . State and prove the analogous result for rational functions \(r \in K = K_{0}(Z_{ij})\) .
\end{enumerate}

\begin{enumerate}
\def\labelenumi{\arabic{enumi})}
\setcounter{enumi}{13}
\tightlist
\item
  Let K be an imperfect field of characteristic p, let a be an element of \(K - K^{p}\) , and let L be the field \(\mathrm{K}(a^{1/p})\) . Consider the K-vector space \(E = L \otimes_{K} L\) ; let u (resp. \(u'\) ) be the K-endomorphism of E given by multiplication by \(a^{1/p} \otimes 1\) (resp. \(1 \otimes a^{1/p}\) ) in the algebra E. Show that u and \(u'\) are semi-simple, that u and \(u'\) commute, and that \(u - u'\) is nonzero nilpotent.
\end{enumerate}

\BourbakiUnnumberedChapter{Historical Note}{Historical Note (Chapters VI and VII)}

(Chapters VI and VII)

(N.B. --- Roman numerals in parentheses refer to the bibliography at the end of this note.)

Elementary arithmetical operations, in particular operations with fractions, inevitably lead to a number of empirical observations about divisibility between integers. But neither the Babylonians (although they were such experts in algebra), nor the Egyptians (despite their acrobatic skill at manipulating fractions) seem to have been acquainted with the general rules governing these properties, and the initiative in this respect is due to the Greeks. Their arithmetical work, a masterly exposition of which can be found in Books VII and IX of Euclid (I), is in no way inferior to their most beautiful discoveries in other branches of Mathematics. The existence of the gcd of two integers is proved right at the beginning of Book VII by the procedure known as the « Euclidean algorithm » *; it serves as the basis for all the subsequent developments (properties of prime numbers, existence and calculation of lcm, etc.), based on arguments not substantially different from those in Chap. VI, § 1 above; and the crowning glory is formed by the two remarkable theorems proving the existence of infinity of prime numbers (Book IX, Prop. 20) and giving a procedure for constructing even perfect numbers (cf.~VII, p.~53, Ex. 24; this procedure in fact gives all the even perfect numbers, as would be proved by Euler). It is only the existence and uniqueness of prime factorisation that is not proved in a general manner; however Euclid does prove explicitly that every integer is divisible by a prime number (Book VII, Prop. 31), as well as the following two propositions (Book IX, Prop. 13 and 14):

« If as many numbers as we like are in a progression, beginning with 1, with constant ratio {[}i.e.~geometric{]}, and if the number following 1 is prime, then the largest number will not be divisible by any number except those appearing in the progression » (in other words, a power \(p^n\) of a prime number \(p\) can be divisible only by powers of \(p\) of exponent \(\leqslant n\) ).

« If a number is the smallest which is divisible by {[}given{]} prime numbers, it will not be divisible by any other prime number with the exception of those initially {[}given as{]} dividing it » (in other words, a product of distinct primes \(p_{1}, \ldots, p_{k}\) has no prime factor other than \(p_{1}, \ldots, p_{k}\) ).

Thus it seems that if Euclid does not state the general theorem, it is only for want of an adequate terminology and notation for arbitrary powers of an integer. *

Although a careful study makes it seem likely that Euclid's text is made up of several successive layers, each corresponding to a stage in the development of Arithmetic **, it seems that this evolution took place entirely between the beginning of the 5th and the middle of the 4th centuries B.C., and one must admire the skill and the logical confidence which it represents: the next comparable progress made in Arithmetic would not occur for two millenia.

It was the so-called « indeterminate » or « Diophantine » problems which were at the root of subsequent developments in Number Theory. The term « Diophantine equations », as it is used today, is not entirely justified historically ; it is generally understood to mean polynomial equations (or systems of equations) with integer coefficients, for which only integer solutions are sought : a problem which is usually impossible if the equations are « determinate », that is have only finitely many (real or complex) solutions, but which on the other hand often admits solutions when there are more unknowns than equations. Now, while Diophantus does indeed seem to have been the first to consider « indeterminate » problems, he looks for integer solutions only in exceptional cases, and is usually content to find a single solution in the rational numbers (II). That was a type of problem which he was usually able to solve by algebraic calculations, where the

* In support of this thesis, one can also remark that the proof of the theorem on perfect numbers is basically just another special case of the unique prime factorisation theorem. Moreover, all the evidence shows that from this period onwards the factorisation of an explicit number into primes was well known and in current use; but no complete proof of the factorisation theorem is to be found before that given by Gauss at the beginning of the Disquisitiones ((VIII) t. I, p.~15).

** Cf. B. L. van der WAERDEN, Die Arithmetik der Pythagoreer, Math. Ann., vol.~CXX (1947-49), p.~127. An example of a passage left over from a previous version is provided by Prop. 21 to 34 of Book IX, which deals with the most elementary properties of divisibility by 2, and no doubt goes back to a period when the general theory of prime numbers had not been developed. Moreover, it is known that the categories Even and Odd played a major rôle in the metaphysical speculations of the early Pythagoreans, to whom it is naturally tempting to attribute this segment.

arithmetical nature of the unknowns was not relevant *; also the theory of divisibility plays only a minor rôle (the word for prime number is used only once ((II), Book V, problem 9, vol.~1, pp.~334-335), and the notion of coprime numbers is invoked only in connection with the theorem which asserts that the quotient of two coprime numbers can be a square only if each of them is a square) **.

The study of integer solutions of indeterminate equations only begins properly with the Chinese and Hindu mathematicians of the early middle ages. The former seem to have been led to considerations of this nature by the practical problems of drawing up calendars (where the determination of the common periods of various cycles of astronomical phenomena constitutes precisely a linear « Diophantine » problem); at any rate they were responsible (certainly between the 4th and the 7th century A.D.) for a rule for solving simultaneous linear congruences (cf.~VI, p.~33, Ex. 25). As for the Hindus, whose Mathematics flourished between the 5th and the 13th centuries, not only did they know how to deal methodically (by applying the Euclidean algorithm) with systems of linear Diophantine equations in arbitrarily many unknowns ***, but they were the first to attack and solve quadratic problems, among them certain special cases of « Fermat's equation » \(Nx^{2} + 1 = y^{2}\) ((III), vol.~II, pp.~87-307).

This is not the place for us to pursue the history of the theory of Diophantine equations of degree \textgreater1, which, by way of the work of Fermat, Euler, Lagrange and Gauss, was to lead to the theory of algebraic integers in the 19th century. As we have already noted (cf.~Historical Note to Chap. II and III), the study of linear systems, which no longer seemed to present problems worthy of interest, was rather neglected during this period: in particular there was no search for general existence conditions for an arbitrary system, nor for a description of the set of solutions. Nevertheless Hermite, in his Number-theoretical research in the 19th

* While Diophantus' work on indeterminate problems is always reduced to problems in a single unknown, via a numerical choice of the other unknowns in such a way as to make a solution to his final equation possible, it seems that the main reason for using this method was his notation, which did not allow him to calculate with several unknowns at once; in any case he keeps track of the numerical substitutions he has made throughout the calculation, and modifies them later if need be, by writing down a compatibility condition for the substituted variables, and solving this auxiliary problem first. In other words, he handles these substituted numerical values like we handle parameters, so much so that what he actually does amounts to finding a rational parametric representation of a given algebraic variety, or of a subvariety of it (cf.~(II bis)).

century, was led to make use of several Lemmas on linear Diophantine equations, notably a « reduced form » for a linear transformation with integer coefficients ((XIII), pp.~164 and 265); finally, after Heger in 1858 had given the existence condition for a system whose rank is equal to the number of equations, H. J. Smith in 1861 defined the invariant factors of an integer matrix, and obtained the general theorem that such a matrix reduces to the « normal form » which we gave in VII, p.~22, Cor. 1 (XVII).

But meanwhile, following its introduction by Gauss (cf.~Historical Notes to Chap. I, II and III), and the important part it played in the subsequent development of Number Theory, the notion of an abelian group was gradually being made precise. In his particularly deep study, presented in the Disquisitiones, of the finite abelian group of classes of quadratic forms of a given discriminant, Gauss soon realised that certain of these groups are not cyclic: « in this case, » he says, « one basis {[}that is to say, one generator{]} cannot suffice, it is necessary to take two or a greater number which, through multiplication and composition *, can produce all the classes » ((VIII), vol.~I, pp.~374-375). It is not certain that Gauss intended these words to describe the decomposition of the group as a direct product of cyclic groups; nevertheless, in the same article of the Disquisitiones, he proves that there exists an element in the group whose order is the lcm of the orders of all the elements --- in other words he obtains the existence of the largest invariant factor of the group ((VIII), vol.~1, p.~373); and on the other hand, the notion of direct product was known to him, for, in a manuscript dating from 1801, but not published during his lifetime, he sketches a general proof of the decomposition of a finite abelian group into a direct product of \(p\) -groups ** ((VIII), vol.~II, p.~226). In any case, in 1868 Schering, the editor of Gauss' collected works, inspired by these results (notably by this manuscript which he had just found), proved (still for the group of classes of quadratic forms) the general decomposition theorem (XVIII) by a method which, as repeated in abstract terms two years later by Kronecker (XX), is essentially that which we have used above (VII, p.~18, Th. 1). As for torsion-free abelian groups, we have already said (cf.~Historical Note to Chap. II and III) how the theory of elliptic functions and abelian integrals, developed by Gauss, Abel and Jacobi, gradually led to attention being paid to their structure; the first and best known example of a decomposition of an infinite group into a direct sum of cyclic groups was given in 1846 by Dirichlet in his paper on the units of an algebraic number field (XI). But it was not until 1879 that the connection between the theory of finitely generated abelian groups and Smith's theorem was recognised and explicitly used by Frobenius and Stickelberger ((XXIII), § 10).

Around the same period the theory of similarity of matrices (with real or complex entries) was also reaching completion. The notion of an eigenvalue of a linear transformation appeared explicitly in the theory of systems of linear differential equations with constant coefficients, applied by Lagrange (VIa) to the theory of small oscillations and by Lagrange (VIb) and Laplace (VIIa) to the « secular » perturbations of the planets. It is implicit in many other problems which were also attacked around the middle of the 18th century, such as that of finding the axes of a conic or a quadric (first solved by Euler (Va)), or the study (also developed by Euler (Vb)) of the principal axes of inertia of a solid body (discovered by De Segner in 1755); we know now that it is also involved (in a more disguised form) in the beginnings of the theory of partial differential equations, in particular the equation of a vibrating string. But (leaving this last case aside) the relationship between these various problems was barely recognised before Cauchy (X). Moreover, as most of them involve the use of symmetric matrices, it was principally because of the latter that eigenvalues were initially studied; we will return to this point in more detail in the Historical Notes following the chapters of this treatise devoted to hermitian operators; let us simply note here that, as early as 1826, Cauchy proved the invariance of the eigenvalues of such matrices under similarity, and proved that they are real for a \(3 \times 3\) symmetric matrix (Xa), a result which he generalised three years later (Xb) to arbitrary real symmetric matrices.* The general notion of a projection, introduced by Möbius in 1827, rapidly led to the problem of classifying such transformations (in 2 and 3 dimensions first of all), which is nothing other than the problem of classifying the corresponding matrices up to similarity; but for a long time this question was treated only by the « synthetic » methods which were in vogue in the mid 19th century, and its (in any case rather slow) progress does not seem to have affected the theory of eigenvalues in any way. The same is not true of another geometric question, the classification of « pencils » of conics or quadrics, which from the modern point of view amounts to the study of the elementary divisors of the matrix \(U + \lambda V\) , where \(U\) and \(V\) are two symmetric matrices; it was certainly in this spirit that Sylvester attacked this problem in 1851, carefully examining (in order to find « canonical forms » for the pencil under consideration) what happens to the minors of the matrix \(U + \lambda V\) when a value for \(\lambda\) is substituted which annihilates the determinant (XIV). The purely algebraic aspect of the theory of eigenvalues was progressing simultaneously; so it was that several authors (including Sylvester himself) proved around 1850 that the eigenvalues of \(U^n\) are the \(n\) -th powers of the eigenvalues of \(U\) , while Cayley announced in 1858, in the paper in which he introduced matrix arithmetic (XVI), the « Cayley-Hamilton Theorem » for an arbitrary square matrix *, though he contented himself with a direct proof for \(2 \times 2\) and \(3 \times 3\) matrices. Finally Weierstrass in 1868, using Sylvester's methods, obtained « canonical forms » for a « pencil » \(U + \lambda V\) where this time \(U\) and \(V\) are square matrices, not necessarily symmetric, subject only to the condition that \(\det(U + \lambda V)\) not be identically zero; he deduced from that the definition of the elementary divisors of an arbitrary square (complex) matrix, and proved that they characterise the latter up to similarity (XIX); these results, incidentally, were partially (and apparently independently) recovered by Jordan two years later ** (XXI). Here again it was Frobenius who showed in 1879 that Weierstrass' theorem can easily be deduced from Smith's theorem extended to polynomials ((XXII), § 13); his procedure is the basis of the proof of this theorem which we have given above (VII, p.~35).

We have just referred to the theory of divisibility for polynomials in one variable; the question of division of polynomials must naturally have arisen in the earliest days of algebra, as the inverse operation to multiplication (the latter being known even to Diophantus, at least for polynomials of low degrees); but one can imagine that it is barely possible to attack this problem in a general fashion before a coherent notation for the various powers of the variable has been established. In fact we find very few examples of the « Euclidean » division procedure for polynomials, as we know it, before the middle of the 16th century ***; and S. Stevin (essentially, using exponent notation) seems to have been the first to have had the idea of deducing the extension of the « Euclidean algorithm » to find the gcd of two polynomials ((IV, vol.~I, pp.~54-56). Apart from that, the theory of divisibility had been restricted to the rational integers until the mid 18th century. It was Euler in 1770 who began a new chapter in Arithmetic by somewhat rashly extending the notion of divisibility to the integers of a quadratic extension: seeking to determine the divisors of a number of the form \(x^{2} + cy^{2}\) ( \(x, y\) and \(c\) rational integers), he put

\[
x + y \sqrt {- c} = (p + q \sqrt {- c}) (r + s \sqrt {- c}) (p, q, r, s \text { rational   integers })
\]

and, taking norms of each side, he had no hesitation in asserting that all the divisors of \(x^{2} + cy^{2}\) are obtained in this way as \(p^{2} + cq^{2}\) (Vc). In other words,

Euler argued as if the ring \(Z[\sqrt{-c}]\) were a principal ideal domain; a little later he used an analogous argument to apply the method of « infinite descent » to the equation \(x^{3} + y^{3} = z^{3}\) (he reduced the problem to that of finding a cube root of \(p^{2} + 3q^{2}\) , which he does by putting \(p + q\sqrt{-3} = (r + s\sqrt{-3})^{3}\) ). But Lagrange proved as early as 1773 (VIc) that the divisors of numbers of the form \(x^{2} + cy^{2}\) are not all of this form, the first example of the fundamental difficulty which was to appear much more clearly in the studies of Gauss and his successors into divisibility in cyclotomic fields *; it is not possible, in general, to extend the essential properties of divisibility of rational integers, such as the existence of gcd's and the uniqueness of prime factorisation, directly to these fields. This is not the place to describe how Kummer, for cyclotomic fields (XII)**, and then Dedekind and Kronecker for arbitrary algebraic number fields, succeeded in overcoming this formidable obstacle by the invention of ideal theory, one of the most decisive advances of modern algebra. But Dedekind, ever curious about the foundations of various Mathematical theories, was not content with this success; and by analysing the mechanism of the divisibility relations he laid the foundations of the theory of lattice ordered groups, in a paper (unknown to his contemporaries, and lost in obscurity for 30 years) which is without doubt one of the earliest works of axiomatic Algebra (XXIV); modulo notation, his work was very close to the modern form of this theory, as we have presented it in Chap. VI, § 1.

From the middle of the 18th century, the order of the day was a search for a « fundamental theorem of algebra » (cf.~Historical Note to Chap. IV and V). We do not need to mention here the attempt of d'Alembert, who initiated the series of proofs which use infinitesimal calculus (cf.~Gen.~Top., Historical Note to Chap. VIII). But Euler attacked the problem in a completely different manner in 1749

* Gauss seems to have hoped at one time that the ring of integers in the field of \(n\) -th roots of unity would be a principal ideal domain; in a manuscript unpublished in his lifetime ((VIII), vol.~II, pp.~387-397), he proves the existence of a Euclidean division process in the field of cube roots of unity, and gives some indications of an analogous process in the field of 5th roots of unity; he uses these results to prove by an « infinite descent » argument more correct than Euler's that the equation \(x^3 + y^3 = z^3\) has no solution in the field of cube roots of unity, indicates that one can extend the method to the equation \(x^5 + y^5 = z^5\) , but stops short of the equation \(x^7 + y^7 = z^7\) , saying that here it is impossible to reject a priori the case where \(x, y\) and \(z\) are not divisible by 7.

(Vd): for every polynomial f with real coefficients, he attempted to prove the existence of a factorisation \(f = f_{1}f_{2}\) into two (non constant) polynomials with real coefficients, which would give him a proof of the « fundamental theorem » by induction on the degree of f.~It even suffices, as he noticed, to stop at the first odd degree factor, and consequently all the difficulties are reduced to the consideration of the case where the degree n of f is even. Euler then restricts himself to the case when the desired factors each have degree n/2, and he indicates that by a suitable process of elimination one can express the unknown coefficients of \(f_{1}\) and \(f_{2}\) as rational functions of a root of an equation with real coefficients, whose extreme terms have opposite signs and which consequently has at least one real root. But Euler's proof was only a sketch, and skipped over a number of essential points; it was not till 1772 that Lagrange succeeded in resolving the difficulties raised by this proof (VId) by means of an extremely long and minute analysis, in which he demonstrated a remarkable virtuosity in the use of the « Galois methods » which he had recently created (cf.~Historical Note to Chap. IV and V).

All the same, Lagrange, like Euler and all his contemporaries, had no hesitation in arguing formally within a « field of roots » of a polynomial (that is, in his language, to consider « imaginary roots » of this polynomial); the Mathematics of his period had provided no justification for this type of argument. Gauss, who was from the outset resolutely hostile to the unrestricted formalism of the 18th century, came out strongly against this abuse in his dissertation ((VIII), vol.~III, p.~3). But it would have been unlike him not to have sensed that it was a case of a superficially faulty presentation of an argument which was intrinsically correct. We find him also, a few years later ((VIII), vol.~III, p.~33; cf.~also (VIIIBis)), taking up a simpler form of Euler's argument, which had been suggested in 1759 by Foncenex (who, however, had been unable to use it to any advantage), to obtain a new proof of the « fundamental theorem » in which he carefully avoids all mention of « imaginary » roots: the latter being replaced by skilful adjunction and specialisation of indeterminates. It is essentially this proof of Gauss that we have given in the text (VI, p.~26, Th. 3), with simplifications made possible by the use of algebraic extensions.

The rôle of Topology in the « fundamental theorem » was thus reduced to the single theorem that a polynomial with real coefficients cannot change sign in an interval without having a zero (Bolzano's Theorem for polynomials). This theorem is also at the root of all the criteria for separation of the real roots of a polynomial (with real coefficients), which was one of the favourite Algebraic topics in the 19th century *. In the course of this research it becomes obvious that it is the order structure of R, much more than its topology, which plays the essential rôle * ; for example Bolzano's Theorem for polynomials remains true in the field of all real algebraic numbers. This train of ideas resulted in the abstract theory of ordered fields, created by E. Artin and O. Schreier (XXV) ; one of its most remarkable results is certainly the discovery that the existence of an order relation on a field is related to purely algebraic properties of the field. This is the theory presented in § 2 of Chap. VI.

\BourbakiUnnumberedChapter{Bibliography}{Bibliography}

\begin{enumerate}
\def\labelenumi{(\Roman{enumi})}
\tightlist
\item
  Euclidis Elementa, 5 vol., ed.~J. L. Heiberg, Lipsiae (Teubner), 1883-88.
\end{enumerate}

(I bis) T. L. HEATH, The thirteen books of Euclid's Elements\ldots, 3 vol., Cambridge, 1908.

\begin{enumerate}
\def\labelenumi{(\Roman{enumi})}
\setcounter{enumi}{1}
\tightlist
\item
  Diophanti Alexandrini Opera Omnia\ldots, 2 vol., ed.~P. Tannery, Lipsiae (Teubner), 1893-95.
\end{enumerate}

(II bis) T. L. HEATH, Diophantus of Alexandria, 2 \(^{nd}\) ed., Cambridge, 1910.

\begin{enumerate}
\def\labelenumi{(\Roman{enumi})}
\setcounter{enumi}{2}
\item
  B. DATTA and A. N. SINGH, History of Hindu Mathematics, 2 vol., Lahore (Motilal Banarsi Das), 1935-38.
\item
  S. STEVIN, Les œuvres mathématiques\ldots, ed.~A. Girard, Leyde (Elsevier), 1634, vol.~I.
\item
  L. EULER: a) Introductio in Analysin Infinitorum (Opera Omnia, (1), t. IX, Zürich-Leipzig-Berlin (O. Füssli and B. G. Teubner), 1945, p.~384); b) Theoria motus corporum solidorum seu rigidorum (Opera Omnia (2), t. III, Zürich-Leipzig-Berlin (O. Füssli and B. G. Teubner), 1948, p.~200-201); c) Vollständige Anleitung zur Algebra (Opera Omnia (1), t. I, Leipzig-Berlin (Teubner), 1911, p.~422); d) Recherches sur les racines imaginaires des équations (Opera Omnia (1), t. VI, Leipzig-Berlin (Teubner), 1921, p.~78).
\item
  J.-L. LAGRANGE, Œuvres, Paris (Gauthier-Villars), 1867-1892: a) Solutions de divers problèmes de Calcul intégral, t. I, p.~520; b) Recherches sur les équations séculaires du mouvement des nœuds, t. VI, p.~655-666; c) Recherches d'arithmétique, t. III, p.~695-795; d) Sur la forme des racines imaginaires des équations, t. III, p.~479; e) Nouvelle solution du problème de rotation d'un corps quelconque qui n'est animé par aucune force accélératrice, t. III, p.~579-616.
\item
  P. S. LAPLACE : a) Mémoire sur les solutions particulières des équations différentielles et sur les inégalités séculaires des planètes (Œuvres, t. VIII, Paris (Gauthier-Villars), 1891, p.~325-366) ; b) Mémoire sur les inégalités séculaires des planètes et des satellites (Œuvres, t. XI, Paris (Gauthier-Villars), 1895, p.~49-92).
\item
  C. F. GAUSS, Werke, t. I (Göttingen, 1870), t. II (ibid., 1876) et t. III (ibid., 1876).
\end{enumerate}

(VIII bis) Die vier Gauss'schen Beweise für die Zerlegung ganzer algebraischer Functionen in reelle Factoren ersten oder zweiten Grades (Ostwald's Klassiker, n° 14, Leipzig (Teubner), 1904).

\begin{enumerate}
\def\labelenumi{(\Roman{enumi})}
\setcounter{enumi}{8}
\item
  N. H. ABEL, Œuvres, t. I, ed.~Sylow and Lie, Christiania, 1881.
\item
  A. L. CAUCHY : a) Leçons sur les applications du Calcul infinitésimal à la Géométrie (Œuvres complètes (2), t. V, Paris (Gauthier-Villars), 1903, p.~248) ; b) Sur l'équation à l'aide de laquelle on détermine les inégalités séculaires des planètes (Œuvres complètes (2), t. IX, Paris (Gauthier-Villars), 1891, p.~174).
\item
  P. G. LEJEUNE-DIRICHLET, Werke, t. I, Berlin (G. Reimer), 1889, p.~619-644.
\item
  E. KUMMER, Zur Theorie der complexen Zahlen, J. de Crelle, t. XLIII (1847), p.~319 (Collected papers, vol.~I, Heidelberg (Springer V.), 1975, p.~203).
\item
  Ch. HERMITE, Œuvres, t. I, Paris (Gauthier-Villars), 1905.
\item
  J. J. SYLVESTER, Collected Mathematical Papers, vol.~I, Cambridge, 1904: An enumeration of the contacts of lines and surfaces of the second order, p.~219 (= Phil. Mag., 1851).
\item
  W. R. HAMILTON, Lectures on Quaternions, Dublin, 1853.
\item
  A. CAYLEY, Collected Mathematical Papers, Cambridge, 1889-1898: A memoir on the theory of matrices, t. II, p.~475-496 (= Phil. Trans., 1858).
\item
  H. J. SMITH, Collected Mathematical Papers, vol.~I, Oxford, 1894; On systems of linear indeterminate equations and congruences, p.~367 (Phil. Trans., 1861).
\item
  E. SCHERING, Die fundamental Classen der zusammensetzbaren arithmetischen Formen, Abh. Ges. Göttingen, t. XIV (1868-69), p.~13.
\item
  K. WEIERSTRASS, Mathematische Werke, Bd. II, Berlin (Mayer und Müller), 1895: Zur Theorie der bilinearen und quadratischen Formen, p.~19.
\item
  L. KRONECKER, Auseinandersetzungen einiger Eigenschaften der Klassenanzahl idealer complexer Zahlen, Monats. Abhandl. Berlin (1870), p.~881 (= Werke, t. I, Leipzig (Teubner), 1895, p.~273).
\item
  C. JORDAN, Traité des substitutions et des équations algébriques. Paris (Gauthier-Villars), 1870, p.~114-125.
\item
  G. FROBENIUS, Theorie der linearen Formen mit ganzen Coefficienten, Gesammelte Abhandlungen, vol.~I, Heidelberg (Springer V.), 1968, p.~482 (= J. de Crelle, 1879).
\item
  G. FROBENIUS und L. STICKELBERGER, Ueber Gruppen von vertauschbaren Elementen, J. de Crelle, t. LXXXVI (1879), p.~217 (= Frobenius, Ges. Abh., vol.~I, p.~545).
\item
  R. DEDEKIND, Gesammelte mathematische Werke, t. II, Braunschweig (Vieweg), 1932: Ueber Zerlegungen von Zahlen durch ihre grössten gemeinsamen Teiler, p.~103.
\item
  \begin{enumerate}
  \def\labelenumii{\alph{enumii})}
  \tightlist
  \item
    E. ARTIN und O. SCHREIER, Algebraische Konstruktion reeler Körper, Abh. Math. Sem. Univ. Hamburg, t. V (1927), p.~83; b) E. ARTIN, Ueber die Zerlegung definite Funktionen in Quadrate (ibid., p.~100); c) E. ARTIN und O. SCHREIER, Eine Kennzeichnung der reell abgeschlossenen Körper (ibid., p.~225).
  \end{enumerate}
\end{enumerate}

\BourbakiUnnumberedChapter{Index of notation}{Index of notation}

\[
\mathbf {A} \left[ \left(\mathbf {X} _ {i}\right) _ {i \in I} \right], \mathbf {A} \left[ \mathbf {X} _ {1}, \dots , \mathbf {X} _ {n} \right], \mathbf {X} ^ {\nu}: \mathrm{IV}, \text { p. } 1.
\]

\[
\deg u: \mathrm{IV}, \text { p. } 3.
\]

\[
\mathbf {f} \circ \mathbf {g}: \mathrm{IV}, \text { p. } 4.
\]

\[
\mathrm{D} _ {i} \mathrm{P}, \mathrm{D} _ {\mathrm{X} _ {i}} \mathrm{P}, \frac {\partial \mathrm{P}}{\partial \mathrm{X} _ {i}}, \mathrm{P} _ {\mathrm{X} _ {i}} ^ {\prime}, \mathrm{D} ^ {\nu}, \mathrm{DP}, \frac {d \mathrm{P}}{d \mathrm{X}}, \mathrm{P} ^ {\prime}: \mathrm{IV}, \text { p. } 6.
\]

\[
\Delta^ {\nu}: \mathrm{IV}, \text { p. } 7.
\]

\[
\mathbf {K} ((\mathbf {X} _ {i}) _ {i \in I}), \deg r: \mathrm{IV}, \text { p. } 19 \text { and } 20.
\]

\[
\mathrm{D} _ {i} f, \mathrm{D} _ {\mathbf {X} _ {i}} f, \frac {\partial f}{\partial \mathbf {X} _ {i}}, f _ {\mathbf {X} _ {i}} ^ {\prime}, \mathrm{D} f, \frac {d f}{d \mathbf {X}}, f ^ {\prime}: \mathrm{IV}, \text {p.} 23.
\]

\[
\mathbf {A} \left[ \left[ \left(\mathbf {X} _ {i}\right) _ {i \in I} \right] \right], \mathbf {A} \left[ [ \mathrm{I} ] \right]: \mathrm{IV}, \text { p. } 24.
\]

\[
\omega : \mathrm{IV}, \text { p. } 25.
\]

\[
u (\mathbf {x}), u ((x _ {i}) _ {i \in I}), u (x _ {1}, \dots , x _ {n}): \mathrm{IV}, \text { p. } 29.
\]

\[
\Delta^ {\nu}, \mathrm{D} ^ {\nu}, \mathrm{D} _ {i}: \mathrm{IV}, \text { p. } 31 \text { and } 32.
\]

\[
\mathrm{D} _ {i} u, \mathrm{D} _ {\mathbf {X} _ {i}} u, \frac {\partial u}{\partial \mathbf {X} _ {i}}, u _ {\mathbf {X} _ {i}} ^ {\prime}, \mathrm{D} u, \frac {d u}{d X}: \mathrm{IV}, \mathrm{p.32}.
\]

\[
\mathrm{A} \{\mathrm{I} \}, \mathbf {f} \circ \mathbf {g}: \mathrm{IV}, \text { p. } 29 \text { and } 36.
\]

\[
\mathrm{T} _ {\mathbf {f}}: \mathrm{IV}, \text { p. } 36.
\]

\[
\exp \mathrm{X}, e ^ {\mathrm{X}}, e (\mathrm{X}), l (\mathrm{X}): \text { IV }, \text { p. } 39 \text { and } 40.
\]

\[
\exp f, \log g: \mathrm{IV}, \text { p. } 40.
\]

\[
\mathbf {M} ^ {\mathrm{H}}, \operatorname{Tr} _ {\mathrm{H} / \mathrm{G}}: \mathrm{IV}, \text { p. } 41 \text { and } 42.
\]

\[
\mathbf {T S} ^ {n} (\mathbf {M}), \mathbf {T S} (\mathbf {M}): \mathrm{IV}, \text { p. } 42.
\]

\[
\mathfrak {S} _ {p \mid q}, \mathfrak {S} _ {p, q}, \mathfrak {S} _ {p _ {1} \mid \dots \mid p _ {n}}: \mathrm{IV}, \text {   p.   } 43.
\]

\[
\gamma_ {k} (x), x \in \mathrm{M}: \mathrm{IV}, \text { p. } 45.
\]

\[
\varphi_ {\mathbf {M}}, \psi_ {\mathbf {M}}: \mathrm{IV}, \text { p. } 52.
\]

\[
\operatorname{Pol} _ {\mathbf {A}} ^ {q} (\mathbf {M}, \mathbf {N}), \operatorname{Pol} ^ {q} (\mathbf {M}, \mathbf {N}): \mathrm{IV}, \text { p. } 55.
\]

\[
\operatorname{Map} (\mathbf {M}, \mathbf {N}), \operatorname{Pol} _ {\mathbf {A}} (\mathbf {M}, \mathbf {N}), \operatorname{Pol} (\mathbf {M}, \mathbf {N}): \mathrm{IV}, \text { p. } 57.
\]

\[
s _ {k}, s _ {k, n}, \mathrm{A} [ \mathrm{X} _ {1}, \dots , \mathrm{X} _ {n} ] ^ {\text { sym }}: \mathrm{IV}, \text { p. } 61.
\]

\[
\mathrm{S} (\alpha), \mathrm{M} (\alpha): \mathrm{IV}, \text { p. } 65 \text { and } 66.
\]

\[
s _ {k}, \mathrm{A} [ [ \mathrm{X} ] ] ^ {\text { sym }}: \mathrm{IV}, \text { p. } 67 \text { and } 68.
\]

\[
\mathfrak {B} _ {k}: \mathrm{IV}, \text { p. } 70.
\]

\[
M (f, g, p, q), \operatorname{res} _ {p, q} (f, g), \operatorname{res} (f, g): \mathrm{IV}, \text {   p.   } 76.
\]

\[
\operatorname{dis} (f), f \text {   monic   polynomial }: \mathrm{IV}, \mathrm{p.~81}.
\]

\[
\operatorname{dis} _ {m} (f), f \text { polynomial   of   degree } \leqslant m: \mathrm{IV}, \text { p. } 83.
\]

\[
\Gamma (\mathrm{E}), \Gamma_ {p} (\mathrm{E}), \gamma_ {p}, \Gamma (h): \mathrm{IV}, \text { p. } 92, \text { ex. } 2.
\]

\[
\mathbf {Q}, \mathbf {F} _ {p}: \mathbf {V}, \text { p. } 1.
\]

\[
\mathbf {S} ^ {p ^ {f}}: \mathbf {V}, \mathrm{p.} 4.
\]

\(U_{m,\alpha}\) : VII, p.~35.

\[
\mathrm{K} [ \mathrm{S} ]: \mathrm{V}, \text { p. } 4.
\]

\[
\mathbf {S} ^ {p ^ {- f}}, \mathbf {A} ^ {p ^ {- \infty}}: \mathbf {V}, \text { p. } 5 \text { and } 6.
\]

\[
[ \mathbf {A}: \mathbf {K} ]: \mathbf {V}, \text { p. } 10.
\]

\[
\mathrm{K} (x _ {i}), \mathrm{K} (\mathrm{x}), \mathrm{K} (x _ {1}, \dots , x _ {n}): \mathrm{V}, \text { p. } 10.
\]

\[
h (\mathrm{L}), [ \mathrm{A}: \mathrm{K} ] _ {s}, \mathscr {H} (\mathrm{A}): \mathrm{V}, \text { p. } 31.
\]

\[
\mathrm{E} _ {s}: \mathrm{V}, \text { p. } 44.
\]

\[
[ \mathrm{E}: \mathrm{K} ] _ {s}, [ \mathrm{E}: \mathrm{K} ] _ {i}: \mathrm{V}, \text { p. } 31 \text { and } 46.
\]

\[
\mathrm{N} _ {\mathrm{A} / \mathrm{K}} (x), \operatorname{Tr} _ {\mathrm{A} / \mathrm{K}} (x), \mathrm{D} _ {\mathrm{A} / \mathrm{K}} (x _ {1}, \dots , x _ {n}): \mathrm{V}, \text { p. } 47.
\]

\[
\operatorname{Gal} (\mathrm{N} / \mathrm{K}): \mathrm{V}, \text { p. } 58.
\]

\[
k (\Delta), g (E): V, p. 67.
\]

\[
\mathbf {K} _ {ab}: \mathbf {V}, \text { p. } 77.
\]

\[
\mu_ {m} (\mathbf {K}), \mu_ {\infty} (\mathbf {K}), \mathbf {Z} [ 1 / p ]: \mathbf {V}, \text {   p.   } 78.
\]

\[
\mu_ {l\infty} (\mathbf {K}): \mathbf {V}, \text { p. } 79.
\]

\[
\varphi (n): \mathrm{V}, \text { p. } 79.
\]

\[
\mathrm{R} _ {n} (\mathrm{K}), \Phi_ {n}, \chi_ {n}: \mathrm{V}, \text { p. } 81.
\]

\[
\mathrm{K} \left(\mathrm{A} ^ {1 / n}\right), \langle \sigma , a \rangle : \mathrm{V}, \text { p. } 88.
\]

\[
\wp , \mathrm{K} (\wp^ {- 1} (\mathrm{A})), [ \sigma , a \rangle : \mathrm{V}, \text { p. } 91.
\]

\[
\mathbf {F} _ {q} (\Omega), \mathbf {F} _ {q}: \mathbf {V}, \text { p. } 95.
\]

\[
\mathbf {Z} _ {l}, \hat {\mathbf {Z}}: \mathbf {V}, \mathrm{p.} 96.
\]

\[
\sigma_ {q}: \mathrm{V}, \text { p. } 97.
\]

\[
\varphi_ {n}: \mathrm{V}, \text { p. } 97.
\]

deg.tr \(_{K}\) E : V, p.~110.

\[
f ^ {\Delta}: \mathbf {V}, \text { p. } 127.
\]

\(\chi(f), f\) a linear mapping: V, p.~132.

\[
x \mid y, x \mid y: \mathrm{VI}, \text { p. } 5.
\]

\[
(x): \mathrm{VI}, \text { p. } 6.
\]

\[
x \equiv x ^ {\prime} (\mathrm{mod} y): \mathrm{VI}, \text { p. } 6.
\]

\(\sup_{F}(\mathbf{x}_{i})\) (F a subset of an ordered set E): VI, p.~8.

\(\gcd (x_{i}), \operatorname{lcm}(x_{i}) : \mathrm{VI}, \text{p. 8.}\)

\(x^{+}, x^{-}, |x|\) (x an element of a lattice ordered group): VI, p.~12.

\[
\operatorname{sgn} (x): \mathrm{VI}, \text { p. } 20.
\]

\(\sqrt{a}\) ( \(a\) an element \(\geqslant 0\) of an ordered field) VI, p.~24.

\(|z|\) (z an element of \(K(i)\) , where \(K\) is an ordered field and \(i^2 = -1\) ): VI, p.~27.

\(\mathbf{GL}^{+}\) (E) (E an oriented vector space): VI, p.~29.

\(M(\alpha), M_{\alpha}\) (M a module over A, \(\alpha \in A\) ): VII, p.~7.

\((\mathbf{Z} / n\mathbf{Z})^{*}, \mathrm{U}(p^{n})\) : VII, p.~12.

\(c_{\mathsf{L}}(x)\) , L a free module over a principal ideal domain A, \(x\in \mathbf{L}\) : VII, p.~16.

\(m(0), m(p^n)\) ( \(p\) irreducible, \(n\) an integer \(\geqslant 1\) ): VII, p.~24.

D(M) (M a module over a principal ideal domain A, \(N^{0}\) : VII, p.~25 to 27.

\[
c _ {\mathbf {M}}: \mathbf {M} \rightarrow \mathbf {D} (\mathbf {D} (\mathbf {M})): \text { VII,   p. } 26.
\]

\(\mathbf{M}_u\) (M a module, \(u\) an endomorphism of \(\mathbf{M}\) ): VII, p.~28.

\(V_{\alpha}\) (V a vector space, \(\alpha \in \mathbf{K}\) ): VII, p.~30.

\(\chi_{u}, \chi_{U}\) : VII, p.~30.

\(u_{s}, u_{n}\) (u an endomorphism): VII, p.~43 and 45.

\(f_{s}, f_{u}\) ( \(f\) an automorphism): VII, p.~46.

\BourbakiUnnumberedChapter{Index of terminology}{Index of terminology}

Abelian closure : V, p.~77.

Abelian extension : V, p.~77.

Absolutely semi-simple endomorphism : VII, p.~42.

Addition of inequalities : VI, p.~2.

Adjunction (extension obtained by): V, p.~11.

Algebraic closure : V, p.~22.

Algebraic element in an algebra : V, p.~16.

Algebraic extension : V, p.~17.

Algebraic number: V, p.~20.

Algebraically closed field : V, p.~20.

Algebraically dependent family, subset : V, p.~106.

Algebraically disjoint extensions : V, p.~112.

Algebraically free family of extensions: V, p.~115.

Algebraically free family, subset : IV, p.~4, V, p.~106.

Archimedean ordered group : VI, p.~35, Ex. 31.

Artin's theorem : V, p.~65.

Artin-Schreier theorem : VI, p.~22.

Artin-Schreier theory : V, p.~91.

Associated elements : VI, p.~5.

Bezout's identity: VII, p.~2.

Budan-Fourier rule : VI, p.~42, Ex. 21.

Characteristic exponent of a field : V, p.~7.

Characteristic of a ring: V, p.~2.

Characteristic submodule : VII, p.~67, Ex. 16.

Chinese remainder theorem: VI, p.~33, Ex. 25.

Closed (field, algebraically): V, p.~20.

Closed (field, separably): V, p.~45.

Closed half-line : VI, p.~28.

Closed (relatively algebraically) field : V, p.~19.

Closure (abelian) of a field : V, p.~77.

Closure (algebraic) of a field : V, p.~22.

Closure (algebraic separable): V, p.~45.

Closure (perfect): V, p.~5.

Closure (relative algebraic): V, p.~19.

Closure (relative p-radical): V, p.~25. Closure (relative separable algebraic): V, p.~43. Coefficients of a formal power series: IV, p.~24. Coefficients of a polynomial: IV, p.~1. Compatible (order relation) with a group, monoid structure: VI, p.~1. Compatible (order relation) with a ring structure: VI, p.~19. Compatible (preorder relation) with a commutative monoid structure: VI, p.~3. Complementary orientation: VI, p.~29. Composite extension: V, p.~12. Composition of series: IV, p.~97, Ex. 15. Conjugacy class in Ω: V, p.~52. Conjugate elements: V, p.~52; Conjugate extensions: V, p.~52. Constant polynomial, term: IV, p.~1. Content: VII, p.~16. Coprime elements: VI, p.~13. Coproduct in TS(M): IV, p.~50. Cubic extension: V, p.~10. Cyclic extension: V, p.~85. Cyclic vector space (for u): VII, p.~29. Cyclotomic extension: V, p.~81. Cyclotomic polynomial: V, p.~81.

Decomposable module : VII, p.~23.

Decomposition algebra (universal) of a polynomial : IV, p.~73.

Decomposition theorem: VI, p.~11.

Dedekind's theorem: V, p.~27.

Degree (inseparable) of an extension: V, p.~46.

Degree of an algebra over a field: V, p.~10.

Degree of an algebraic element: V, p.~16.

Degree (total) of a polynomial: IV, p.~3.

Degree (total) of a rational fraction : IV, p.~20.

Degree (separable) of an extension : V, p.~31.

Derivation (partial) in a p-radical extension of height \(\leqslant 1\) : V, p.~103.

Derivative (partial) of a formal power series: IV, p.~32.

Derivative (partial) of a polynomial : IV, p.~6.

Descartes' rule : VI, p.~42, Ex. 21.

Diagonal, diagonalizable set of endomorphisms : VII, p.~40.

Diagonalizable, diagonalized algebra : V, p.~28.

Direct basis : VI, p.~29.

Direction vector : VI, p.~28.

Direction of a half-line : VI, p.~28.

Discriminant of a monic polynomial : IV, p.~81.

Discriminant of a polynomial : IV, p.~83.

Discriminant of a sequence of elements: V, p.~47.

Disjoint (algebraically) extensions : V, p.~112.

Disjoint (linearly) extensions : V, p.~14.

Divided power in TS(M): IV, p.~45.

Divided power of a series: IV, p.~96, Ex. 13.

Divisibility relation : VI, p.~5.

Division (euclidean): IV, p.~10.

Divisor (greatest common): IV, p.~12, VI, p.~8.

Divisor of an element: VI, p.~5.

Divisors (elementary) of a module: VII, p.~24.

Double root : IV, p.~15.

Eigenspace : VII, p.~30.

Eigenvalue : VII, p.~30.

Elementary divisor of a module: VII, p.~24.

Elementary symmetric polynomial : IV, p.~61.

Etale algebra : V, p.~28.

Euclid's lemma : VI, p.~15.

Euclidean domain : VII, p.~49, Ex. 7.

Euler's identity: VI, p.~8, IV, p.~89, Ex. 6.

Euler-Lagrange theorem: VI, p.~26.

Expansion at the origin of a rational fraction : IV, p.~30, V, p.~39.

Exponential of a formal power series: IV, p.~39.

Exponential type (sequence of): IV, p.~91, Ex. 1.

Extension of a field : V, p.~9.

Extension (principle of) of algebraic identities : IV, p.~18.

Finite extension : V, p.~10.

Finitely generated extension : V, p.~11.

Form of degree n: IV, p.~2.

Formal power series : IV, p.~24.

Formal power series (generalized): IV, p.~38.

Fractions (fields of rational): IV, p.~19.

Frobenius homomorphism : V, p.~4, V, p.~92.

Galois extension : V, p.~57.

Galois group of an extension, of a polynomial : V, p.~58.

Galois theory : V, p.~67.

Gamma algebra of a module: IV, p.~92, Ex. 2.

Gaussian integer : VII, p.~1.

Generated (extension): V, p.~11.

Generated (Galois extension): V, p.~57.

Generated (quasi-Galois extension): V, p.~55.

Generating family of an extension : V, p.~11.

Greatest common divisor (GCD) of elements : VI, p.~8.

Greatest common divisor of polynomials : IV, p.~12.

Greatest common divisor of principal ideals: VII, p.~3.

Half-line (closed, open): VI, p.~28.

Half-lines (opposite): VI, p.~28.

Hankel determinants : IV, p.~90, Ex. 1.

Height of a \(p\) -radical element: V, p.~24.

Height of a \(p\) -radical extension: V, p.~26.

Height of an element of \(\mathbf{K}(\mathbf{T})\) : V, p.~148, Ex. 11.

Hermite interpolation formula : IV, p.~88, Ex. 13.

Hermite reduced form : VII, p.~68, Ex. 21.

Homogeneous component of a formal power series : IV, p.~25.

Imperfect field : V, p.~7.

Imperfection (degree of): V, p.~170, Ex. 1.

Indecomposable module : VII, p.~23.

Indeterminate : IV, p.~1.

Index of a linear mapping : V, p.~132.

Indicator (Euler's): V, p.~79.

Indivisible element : VII, p.~17.

Inflation homomorphism : V, p.~70.

Inseparable degree : V, p.~46.

Integral ideal: VI, p.~6.

Invariant factors of a linear mapping : VII, p.~22.

Invariant factors of a module: VII, p.~20.

Invariant factors of a submodule : VII, p.~19.

Invariants (similarity) of an endomorphism : VII, p.~32.

Irreducible element : VI, p.~17.

Irreducible polynomial : IV, p.~13.

Isobaric polynomial : IV, p.~3.

Jordan decomposition : VII, p.~43.

Jordan decomposition (multiplicative): VII, p.~46.

Jordan matrix : VII, p.~35.

Kummer extension : V, p.~162, Ex. 7.

Kummer theory : V, p.~88.

Lagrange interpolation formula : IV, p.~16.

Law (polynomial): IV, p.~94, Ex. 9.

Leading coefficient : IV, p.~3.

Least common multiple (LCM) of elements : VI, p.~8.

Least common multiple of principal ideals: VII, p.~3.

Lefschetz's principle : V, p.~117.

Lexicographic product of ordered groups : VI, p.~7.

Linear topology : IV, p.~28.

Linearly disjoint extensions : V, p.~13.

Linearly topologized algebra : IV, p.~28.

Logarithm of a formal power series : IV, p.~40. Logarithmic derivative : IV, p.~41. Lüroth's theorem : V, p.~149, Ex. 11.

Mac Laner's criterion : V, p.~43, V, p.~123. Maximal ordered field : VI, p.~25. Minimal polynomial of an element : V, p.~16. Minimal polynomial of an endomorphism : VII, p.~29. Minkowski's lemma : VII, p.~50, Ex. 9. Monic polynomial : IV, p.~3. Monogenous extension : V, p.~11. Monomials, IV, p.~1. Multiple factor (polynomial without) : IV, p.~14. Multiple of an element : VI, p.~5. Multiple root : IV, p.~15. Multiplicity (geometric) of an eigenvalue : VII, p.~30. Multiplicity of a root : IV, p.~15. Multiplicity of an elementary divisor : VII, p.~24.

Negative element : VI, p.~4.

Negative \(n\) -vector: VI, p.~29.

Negative basis : VI, p.~29.

Newton polygon : V, p.~150, Ex. 2.

Newton's relations: IV, p.~70, IV, p.~75.

Nilpotent component : VII, p.~45.

Normal basis, normal basis theorem : V, p.~73.

Normal extension : V, p.~53.

Norm : V, p.~47.

Open half-line : VI, p.~28.

Opposite half-lines : VI, p.~28.

Order of a (generalized) formal power series: IV, p.~25, IV, p.~38.

Order of a root: IV, p.~15.

Orderable field : VI, p.~39, Ex. 8.

Ordered extension : VI, p.~21.

Ordered field : VI, p.~20.

Ordered group, monoid : VI, p.~1.

Ordered ring : VI, p.~19.

Orientation of a vector space : VI, p.~29.

Oriented affine space, vector space : VI, p.~29.

\(p\) -adic integers (ring of): V, p.~96.

\(p\) -radical closure: V, p.~25.

\(p\) -radical element: V, p.~24.

\(p\) -radical extension: IV, p.~19.

\(p\) -torsion group: VII, p.~10.

\(\pi\) -primary component: VII, p.~8.

\(\pi\) -primary module: VII, p.~7.

Perfect closure of a ring: V, p.~5.

Perfect field : V, p.~7.

Perfect ring of characteristic p : V, p.~5.

Polynomial, polynomial algebra : IV, p.~1.

Polynomial function : IV, p.~4.

Polynomial law : IV, p.~94, Ex. 9.

Polynomial mapping : IV, p.~57.

Polynomial mapping (homogeneous): IV, p.~55.

Positive basis : VI, p.~29.

Positive element : VI, p.~4.

Positive \(n\) -vector: VI, p.~29.

Power series (generalized formal): IV, p.~24, IV, p.~38.

Preordered group, monoid: VI, p.~3.

Primary extension: V, p.~139.

Prime field, subfield : V, p.~1.

Primitive element (theorem of the): V, p.~40.

Primitive root of unity : V, p.~81.

Principal fractional ideal: VI, p.~6.

Principal ideal domain : VII, p.~1.

Principal ideal ring: VII, p.~49, Ex. 6.

Product (lexicographic) of ordered groups: VI, p.~7.

Product of ordered groups : VI, p.~7.

Product orientation : VI, p.~29.

Puiseux's theorem: V, p.~150, Ex. 2.

Pure basis : V, p.~106.

Pure extension : V, p.~106.

Pure submodule : VII, p.~55, Ex. 7.

Pythagorean field : VI, p.~39, Ex. 8.

Quadratic extension : V, p.~10.

Quadratic reciprocity (theorem of): V, p.~166, Ex. 23.

Quasi-Galois extension : V, p.~53.

Quotient in euclidean division : IV, p.~11.

Quotient orientation : VI, p.~29.

Rational fraction : IV, p.~19.

Rational function : IV, p.~21.

Reduced ring : V, p.~34.

Regular algebra : V, p.~140.

Regular extension : V, p.~141.

Relative algebraic closure of a field in an extension: V, p.~19.

Relatively algebraically closed field in an extension field : V, p.~19. Relatively prime polynomials : IV, p.~12. Remainder in euclidean division : IV, p.~11. Representatives (system of) of irreducible elements : VII, p.~3. Restriction homomorphism : V, p.~60, V, p.~70. Resultant of two polynomials : IV, p.~76. Root of a polynomial : IV, p.~14. Root of unity : V, p.~78. Rule of signs : VI, p.~19.

Semi-simple (absolutely) component : VII, p.~45. Semi-simple (absolutely) endomorphism : VII, p.~42. Semi-simple endomorphism : VII, p.~41. Semi-simple module : VII, p.~9. Separable algebra : V, p.~119. Separable algebraic closure (relative) : V, p.~44. Separable algebraic extension : V, p.~36. Separable closure : V, p.~44. Separable degree : V, p.~31. Separable element : V, p.~39. Separable extension : V, p.~121. Separable polynomial : V, p.~37. Separably closed field : V, p.~45. Separating transcendence basis : V, p.~136. Series on E with values in F : IV, p.~96, Ex. 11. Set (Γ-) : V, p.~75. Sign of an element : VI, p.~20. Similar endomorphisms, matrices : VII, p.~29. Similarity invariants of on endomorphism : VII, p.~32. Simple module : VII, p.~25. Simple root : IV, p.~15. Splitting field, extension : V, p.~21. Subextension : V, p.~9. Subfield : V, p.~1. Substitutable element in a rational fraction : IV, p.~21. Substitution in a formal power series : IV, p.~29. Substitution in a polynomial : IV, p.~4. Substitution in a rational fraction : IV, p.~21. Symmetric (elementary) polynomials : IV, p.~61. Symmetric formal power series : IV, p.~67. Symmetric polynomial : IV, p.~61. Symmetric product of symmetric tensors : IV, p.~43. Symmetric rational fraction : IV, p.~67. Symmetric tensor : IV, p.~42. Symmetrized tensor : IV, p.~43.

Tangent linear mapping : IV, p.~36. Taylor's formula for formal power series : IV, p.~31.

Taylor's formula for polynomials : IV, p.~8. Term, constant term, of a formal power series : IV, p.~24. Term, constant term, of a polynomial : IV, p.~1. Trace : V, p.~47. Transcendence basis : V, p.~109. Transcendence degree : V, p.~110. Transcendental element : V, p.~16. Transcendental extension : V, p.~17. Translation homomorphism : V, p.~70. Triangularisable endomorphism : VII, p.~35. Trivial extension : V, p.~9.

Ulm-Zippin theorem : VII, p.~59, Ex. 14.

Unipotent component : VII, p.~46.

Unipotent endomorphism : VII, p.~46.

Units of a ring: VI, p.~5.

Universal decomposition of a monic polynomial : IV, p.~73.

Weight : IV, p.~3.

Wilson's formula : V, p.~94.

Zero of a polynomial : IV, p.~14.
